\pdfoutput=1
\documentclass[11pt,reqno]{amsart}

\usepackage[T1]{fontenc}
\usepackage[utf8]{inputenc}
\usepackage{lmodern}
\usepackage[final]{microtype}
\usepackage{amsmath,amssymb,amsthm,mathtools}
\usepackage{enumitem}
\usepackage{array}
\usepackage{booktabs}
\usepackage{longtable}
\usepackage{aliascnt}
\usepackage{xcolor}
\usepackage[margin=1.2in]{geometry}
\usepackage[hyphens]{url}
\usepackage{xurl}
\usepackage[colorlinks=true,linkcolor=blue!50!black,citecolor=blue!50!black,
  urlcolor=blue!50!black,
  pdftitle={Morrey's problem for homogeneous isotropic integrands in the plane},
  pdfauthor={The-Anh Vu-Le}]{hyperref}
\usepackage[capitalize,noabbrev]{cleveref}

% Notation of the paper. Section files do not define macros of their own.

% --- numbers, sets
\newcommand{\R}{\mathbb{R}}
\newcommand{\C}{\mathbb{C}}
\newcommand{\M}{\R^{2\times 2}}                    % real 2 x 2 matrices
\newcommand{\abs}[1]{\lvert #1\rvert}
\newcommand{\norm}[1]{\lVert #1\rVert}             % the Euclidean norm of a matrix
\DeclareMathOperator{\re}{Re}
\DeclareMathOperator{\im}{Im}
\DeclareMathOperator{\diag}{diag}
\DeclareMathOperator{\rank}{rank}
\newcommand{\SO}{\mathrm{SO}(2)}
\newcommand{\CO}{\mathrm{CO}(2)}                   % the multiples of orthogonal matrices

% --- the class
\newcommand{\Hyp}{\textup{(H)}}                    % the common hypotheses

% --- Lean declarations
% \lean{name} typesets a declaration name (underscores allowed). Names are upright, are not
% letter-spaced, and break only after a dot, an underscore, a hyphen or a slash.
\renewcommand{\UrlFont}{\ttfamily\upshape}
\Urlmuskip=0mu
\makeatletter
\def\UrlBreaks{\do\.\do\_\do\-\do\/}
\def\UrlBigBreaks{}
\def\UrlNoBreaks{}
\makeatother
\newcommand{\lean}[1]{\textup{\path{#1}}}


% ---------------------------------------------------------------- theorem environments
% All environments share one counter (numbered within sections); the alias counters let cleveref
% recover the type of each label.
\theoremstyle{plain}
% The main theorem of the Introduction is numbered by a letter.
\newtheorem{maintheorem}{Theorem}
\renewcommand{\themaintheorem}{\Alph{maintheorem}}
\crefname{maintheorem}{Theorem}{Theorems}
\newtheorem{theorem}{Theorem}[section]
\newcommand{\shared}[3]{%
  \newaliascnt{#1}{theorem}\newtheorem{#1}[#1]{#2}\aliascntresetthe{#1}\crefname{#1}{#2}{#3}}
\shared{proposition}{Proposition}{Propositions}
\shared{lemma}{Lemma}{Lemmas}
\shared{corollary}{Corollary}{Corollaries}
\theoremstyle{definition}
\shared{definition}{Definition}{Definitions}
\theoremstyle{remark}
\shared{remark}{Remark}{Remarks}

\emergencystretch=1.5em
\vfuzz=2pt % a page may be up to 2pt too high without a warning
\setlength{\textfloatsep}{18pt plus 2pt minus 4pt} % space between a float and the text
\setlength{\intextsep}{18pt plus 2pt minus 4pt}
\setlist[enumerate]{leftmargin=*,itemsep=2pt,topsep=3pt}
\setlist[itemize]{leftmargin=*,itemsep=2pt,topsep=3pt}
\numberwithin{equation}{section}

% The caption of a longtable in the style of the class's table captions.
\setlength{\LTcapwidth}{\textwidth}
\makeatletter
\def\LT@makecaption#1#2#3{%
  \LT@mcol\LT@cols c{\hbox to\z@{\hss\parbox[t]\LTcapwidth{%
    \centering\normalsize #1{\textsc{#2}.\enspace}#3\endgraf\vskip\belowcaptionskip}\hss}}}
\makeatother

% ---------------------------------------------------------------- title
\title[Homogeneous isotropic integrands in the plane]{Morrey's problem for homogeneous
isotropic integrands in the plane}

\author{The-Anh Vu-Le}
% Affiliation and e-mail are the author's to fill in before submission.

\date{November 2026}
\subjclass[2020]{Primary 49J45; Secondary 26B25, 30C62, 68V20}
\keywords{Morrey's problem, rank-one convexity, quasiconvexity, polyconvexity, isotropic
integrands, homogeneous integrands, Burkholder function, truncated determinant, formal proof}

\begin{document}

\begin{abstract}
  Morrey's problem asks whether rank-one convexity implies quasiconvexity. It is open for
  integrands on real $2\times2$ matrices. We prove that the answer is yes for every isotropic,
  positively $p$-homogeneous integrand on these matrices, for every $p>0$. If $p\ge2$, the
  nonnegative integrands of this class are polyconvex, with an explicit polyaffine support at
  every matrix; this answers a question of Guerra and Kristensen. The integrands whose
  quasiconvexity was in question are those that take negative values. For them, the proof
  builds touching quasiconvex minorants by integrating the truncated determinant over scales. It uses the
  quasiconvexity of the truncated determinant, which is claimed in a recent preprint of Agazzi
  et al.\ and which we audited with a Lean formalization. All results are formalized in Lean~4.
\end{abstract}

\maketitle

\setcounter{tocdepth}{1}
{\small\tableofcontents}

\section{Introduction}\label{sec:intro}

Let $\M$ be the space of real $2\times2$ matrices. A function $F\colon\M\to\R$ is
\begin{itemize}
\item \emph{rank-one convex} if $\tau\mapsto F(A+\tau B)$ is convex on $\R$ for every $A$ and
  every $B$ with $\rank B=1$;
\item \emph{quasiconvex} if $F(A)\le\int_{(0,1)^2}F(A+\nabla\varphi)\,dx$ for every $A$ and every
  Lipschitz $\varphi\colon\R^2\to\R^2$ that vanishes outside $(0,1)^2$, whenever the integral
  exists;
\item \emph{polyconvex} if $F(A)=g(A,\det A)$ for a convex $g\colon\M\times\R\to\R$.
\end{itemize}
Polyconvexity implies quasiconvexity, which implies rank-one convexity
\cite{Morrey1952,Ball1977,Dacorogna2008}. Quasiconvexity is the condition under which the
integral $\int F(\nabla u)$ is lower semicontinuous, and it is hard to verify; rank-one convexity
is a condition on lines. \emph{Morrey's problem} asks whether rank-one convexity implies
quasiconvexity. For maps with at least three components the answer is negative, by an example of
\v{S}ver\'ak~\cite{Sverak1992}. A negative answer for maps $\R^m\to\R^2$ with $m\ge3$ is
claimed in the preprint of Cassese~\cite{Cassese2026}, whose proof we audited with a Lean
formalization~\cite{Companion}. For maps $\R^2\to\R^2$, that is, for integrands on $\M$, the
problem is open.

This paper answers it affirmatively for the isotropic, positively homogeneous integrands on
$\M$. We fix an exponent $p>0$ and consider functions $F\colon\M\to\R$ with the following three
properties, which we call the \emph{common hypotheses}~\Hyp.
\begin{enumerate}[label=\textup{(H\arabic*)}]
\item\label{H1} $F$ is \emph{isotropic}: $F(QAR)=F(A)$ for all rotations $Q,R\in\SO$.
\item\label{H2} $F$ is \emph{positively $p$-homogeneous}: $F(\lambda A)=\lambda^pF(A)$ for all
  $\lambda>0$.
\item\label{H3} $F$ is rank-one convex.
\end{enumerate}

\begin{maintheorem}\label{thm:A}
Let $p>0$, and let $F\colon\M\to\R$ satisfy~\Hyp. Then $F$ is quasiconvex.
\end{maintheorem}

For nonnegative integrands, \cref{thm:A} follows from the results of
\cref{sec:intro-large,sec:intro-low,sec:intro-small}, which say more: for $p\ge2$ such an
integrand is polyconvex, with explicit supports, and for $1\le p<2$ it has a touching minorant
built from the Burkholder function at every matrix. For $p>2$, by \cref{thm:iff}\,(a), every
integrand with~\Hyp\ that is not polyconvex takes negative values; the Burkholder function and
the Dacorogna--Marcellini integrands of \cref{sec:intro-related} beyond the polyconvex range are
of this kind. The integrands that take negative values are treated in \cref{sec:signed}.

The best known function with the properties \Hyp\ is the \emph{Burkholder function}, defined
for $1<p<\infty$ and normalized as in~\cite[(1.4)]{GuerraKristensen2022}:
\[
  B_p(A)=\Bigl(\Bigl\lvert1-\frac2p\Bigr\rvert\abs A^2+\det A\Bigr)\abs A^{p-2},\qquad B_p(0)=0,
\]
where $\abs A$ is the operator norm; the same formula defines $B_p$ for every $p>0$, and we
use it also for $p=1$. It comes from Burkholder's sharp inequalities for martingale
transforms~\cite{Burkholder1984,Burkholder1988}, and its quasiconvexity implies Iwaniec's
conjecture~\cite{Iwaniec1982} on the $L^p$ norm of the Beurling--Ahlfors transform
\cite{BaernsteinMontgomerySmith1999,Iwaniec2002,AIPS2012,AIPS2015}. The quasiconvexity of
$B_p$ is claimed in the recent preprint of Agazzi et al.~\cite[Corollary~2]{SharpHI}. The same
preprint claims, as its Theorem~1, the quasiconvexity of the \emph{truncated determinant}. To
write it, identify $\R^2$ with $\C$: a matrix $X$ acts as $\zeta\mapsto z\zeta+w\bar\zeta$ for
two complex numbers $z$, $w$ (\cref{sec:coordinates}), and $\abs X=\abs z+\abs w$,
$\det X=\abs z^2-\abs w^2$. The truncated determinant is
\begin{equation}\label{eq:truncated}
  L(X)=\det X\quad\text{if }\abs X\le1,\qquad L(X)=2\abs z-1\quad\text{if }\abs X>1 .
\end{equation}
Our proofs use Theorem~1, whose proof we audited with a Lean formalization~\cite{IwaniecLean},
and the quasiconvexity of $B_p$ follows from it (\cref{cor:burkholder}).

\subsection{What is proved, and from which hypotheses}\label{sec:intro-table}

\Cref{thm:A} uses only \ref{H1}, \ref{H2} and \ref{H3}. The other results refine it, and some
of them need a fourth hypothesis, the \emph{sign condition}
\begin{enumerate}[label=\textup{(N)}]
\item\label{N} $F\ge0$ on $\M$.
\end{enumerate}
\Cref{tab:summary} lists, for each range of the exponent, what follows from \Hyp\ and what
\ref{N} adds. For $0<p<1$ isotropy is not needed, and for $p\le1$ the sign condition is a
consequence of \Hyp. For $p>2$ it is necessary for polyconvexity.

\begin{table}[ht]
\centering
\caption{The results for each range of the exponent: what follows from \Hyp, and what the sign
  condition \ref{N} adds. Here $J=\diag(1,-1)$. \Cref{thm:A} and \cref{cor:low}\,(a) rest on
  the quasiconvexity of the truncated determinant, claimed as Theorem~1 of~\cite{SharpHI}.}
\label{tab:summary}
\small
\begin{tabular}{@{}l>{\raggedright\arraybackslash}p{0.25\textwidth}%
>{\raggedright\arraybackslash}p{0.33\textwidth}>{\raggedright\arraybackslash}p{0.2\textwidth}@{}}
\toprule
Exponent & From \Hyp & With \ref{N} in addition & Statements \\
\midrule
$0<p<1$ & $F=0$, also without \ref{H1} & --- & \cref{prop:small} \\\addlinespace[3pt]
$p=1$ & $F\ge0$, and $F$ is convex & --- & \cref{cor:low}\,(b), Dacorogna
  \cite{Dacorogna1995} \\\addlinespace[3pt]
$1<p<2$ & $F$ is quasiconvex & at every matrix, a touching minorant $\alpha B_p(A)+\beta B_p(JA)$
  with $\alpha,\beta\ge0$ & \cref{thm:A}; \cref{thm:low}, \cref{cor:low}\,(a)
  \\\addlinespace[3pt]
$p=2$ & $F$ is polyconvex & --- & \cref{thm:iff}\,(b) \\\addlinespace[3pt]
$p>2$ & $F$ is quasiconvex; it is polyconvex if and only if $F\ge0$ & a polyaffine support at
  every matrix, in closed form & \cref{thm:A}, \cref{thm:iff}\,(a);
  \cref{thm:main,thm:closed} \\
\bottomrule
\end{tabular}
\end{table}

\subsection{The exponents \texorpdfstring{$p\ge2$}{p >= 2}}\label{sec:intro-large}

A \emph{polyaffine} function is $P(A)=k+\langle K,A\rangle+c\det A$. Guerra and
Kristensen~\cite{GuerraKristensen2022} studied the nonnegative functions with~\Hyp. Let $\CO$ be
the set of the multiples of orthogonal matrices. They proved that for $p\ge2$ such a function is
polyconvex at the matrices of $\CO$ \cite[Theorem~1.3]{GuerraKristensen2022}, and they
ask~\cite[Question~1.7]{GuerraKristensen2022}: ``When $n=2$, are the integrands satisfying the
assumptions of Theorem~1.3 polyconvex everywhere?'' \Cref{thm:main} says that they are.

\begin{theorem}\label{thm:main}
Let $p\ge2$ and let $F\colon\M\to[0,\infty)$ satisfy~\Hyp. Then for every $A_0\in\M$ there is a
polyaffine function $P$ with $P\le F$ on $\M$ and $P(A_0)=F(A_0)$. In particular $F$ is
polyconvex.
\end{theorem}

At the zero matrix the support is $0$. By isotropy and homogeneity it is then enough to give the
support at the matrices $\diag(1,t)$, $-1\le t\le1$. Put $h(s)=F(\diag(1,s))$ for
$s\in[-1,1]$. Then $h$ is convex, and $F(A)=\abs A^p\,h(\det A/\abs A^2)$ for $A\ne0$, so the
function $h$ of one variable determines $F$.

\begin{theorem}\label{thm:closed}
Let $p$ and $F$ be as in \cref{thm:main}. Let $\abs t<1$, let $H=h(t)>0$, and let $m$ be the slope
of a supporting line of $h$ at $t$, that is, $h(s)\ge H+m(s-t)$ for $s\in[-1,1]$. Put
\[
  L=\frac{(1+t)m-pH}{1-t},\qquad U=\frac{pH+(1-t)m}{1+t},\qquad
  c^\ast=\min\bigl\{U,\max\{L,(p-1)m\}\bigr\}.
\]
Then the polyaffine function
\[
  P_t(A)=c^\ast t-(p-1)H+(pH-tm-c^\ast t)\,a_{11}+(m-c^\ast)\,a_{22}+c^\ast\det A
\]
satisfies $P_t\le F$ on $\M$ and $P_t(\diag(1,t))=F(\diag(1,t))$.
\end{theorem}

The determinant coefficient is $(p-1)m$, the value that an inequality of Burkholder type gives
(\cref{lem:burkholder}), clamped to the interval $[L,U]$. At $t=\pm1$ the support is
$h(\pm1)\bigl(1+\tfrac p2(\pm\det A-1)\bigr)$, and where $h(t)=0$ it is $0$. The support is
computed from $t$, $h(t)$ and the slope of one supporting line. It involves no second derivative
of $h$, and $h$ need not be differentiable.

The sign condition cannot be dropped when $p>2$: the Burkholder function $B_p$ satisfies~\Hyp\
and is not polyconvex \cite[Remark~1.4, Proposition~4.6]{GuerraKristensen2022}. By the next
theorem it is the only condition missing, and for $p=2$ none is needed.

\begin{theorem}\label{thm:iff}
Let $F\colon\M\to\R$ be isotropic and positively $p$-homogeneous.
\begin{enumerate}[label=\textup{(\alph*)}]
\item If $p>2$, then $F$ is polyconvex if and only if it is rank-one convex and nonnegative.
\item If $p=2$, then $F$ is polyconvex if and only if it is rank-one convex.
\end{enumerate}
\end{theorem}

\begin{corollary}\label{cor:equiv}
Let $p\ge2$, and let $F\colon\M\to[0,\infty)$ be Borel, locally bounded, isotropic and
positively $p$-homogeneous. Then $F$ is rank-one convex if and only if it is quasiconvex, if
and only if it is polyconvex.
\end{corollary}

\subsection{The exponents \texorpdfstring{$1\le p<2$}{1 <= p < 2}}\label{sec:intro-low}

For $1<p<2$ polyconvexity fails for functions with~\Hyp: the positive part of $B_p$ is not
polyconvex \cite[Remark~1.4]{GuerraKristensen2022}. The Burkholder function takes the place of
the polyaffine functions. In the conformal coordinates $(z,w)$ of \cref{sec:setting},
$B_p=\tfrac2p(\abs z-(p-1)\abs w)(\abs z+\abs w)^{p-1}$ for $0<p\le2$; in particular
$B_1=2\abs z$. Let $J=\diag(1,-1)$; the function $A\mapsto B_p(JA)$ is the mirror image of
$B_p$, in which $z$ and $w$ change roles.

\begin{theorem}\label{thm:low}
Let $1\le p<2$, and let $F\colon\M\to[0,\infty)$ satisfy~\Hyp. Then for every $A_0\in\M$ there
are $c_0,m\in\R$ with $(2-p)\abs m\le pc_0$ such that the function
\[
  G_{c_0,m}(A)=\abs A^{p-2}\bigl(c_0\abs A^2+m\det A\bigr)\quad(A\ne0),\qquad G_{c_0,m}(0)=0,
\]
satisfies $G_{c_0,m}\le F$ on $\M$ and $G_{c_0,m}(A_0)=F(A_0)$. Moreover
\[
  G_{c_0,m}(A)=\frac{pc_0+(2-p)m}{2(2-p)}\,B_p(A)+\frac{pc_0-(2-p)m}{2(2-p)}\,B_p(JA),
\]
with nonnegative coefficients.
\end{theorem}

The identity holds for every $0<p<2$ and all $c_0,m\in\R$, and its coefficients are
nonnegative exactly when $(2-p)\abs m\le pc_0$. So a nonnegative $F$ with~\Hyp\ is the upper
envelope of nonnegative combinations of the Burkholder function and its mirror image, each
touching $F$ at one matrix. A property of these
two functions passes to $F$ if nonnegative combinations and touching minorants preserve it.
Quasiconvexity and convexity are such properties.

The Burkholder function is quasiconvex for every $p>1$, by the radial integration of the
truncated determinant in \cref{sec:signed} (\cref{prop:radial,cor:burkholder}), and for $p=1$,
$B_1=2\abs z$ is convex.

\begin{corollary}\label{cor:low}
\begin{enumerate}[label=\textup{(\alph*)}]
\item Let $1<p<2$. Every $F\colon\M\to[0,\infty)$ with~\Hyp\ is quasiconvex.
\item Every $F\colon\M\to\R$ with~\Hyp\ for $p=1$ is nonnegative and convex.
\end{enumerate}
\end{corollary}

Part~(b) is a theorem of Dacorogna~\cite{Dacorogna1995}; see also the introduction
of~\cite{DacorognaHaeberly1996}. For the matrices of $\CO$, \cref{thm:low} and part~(a) go back
to Guerra and Kristensen: they proved that $F\ge B_p$ if $F(I)=B_p(I)$, after~\cite{AIPS2015},
and that $F$ is quasiconvex at the matrices of $\CO$ if $B_p$ is quasiconvex
\cite[Theorems~4.2 and~1.2]{GuerraKristensen2022}. Here the minorant touches at every matrix.

Theorem~1 of~\cite{SharpHI} enters the proof of \cref{prop:radial}, and through it those of
part~(a) and of \cref{thm:A}. It is stated there with smooth periodic test maps. Its proof
applies to Lipschitz periodic maps, and for a continuous integrand the periodic test maps and the
Lipschitz maps vanishing outside the unit square define the same notion
\cite[Proposition~5.13]{Dacorogna2008}. The formalization~\cite{IwaniecLean} proves the
quasiconvexity of $L$ in this form, which is the sense used here.

\subsection{The exponents \texorpdfstring{$0<p<1$}{0 < p < 1}}\label{sec:intro-small}

For exponents below $1$ only the zero function remains, and neither isotropy nor a sign
condition is needed.

\begin{proposition}\label{prop:small}
Let $0<p<1$, and let $F\colon\M\to\R$ be positively $p$-homogeneous and rank-one convex. Then
$F=0$.
\end{proposition}

For a nonnegative $F$ with~\Hyp, \cref{thm:A} follows: $F=0$ for $0<p<1$, $F$ is convex for
$p=1$, quasiconvex for $1<p<2$ by \cref{cor:low}\,(a), and polyconvex for $p\ge2$ by
\cref{thm:main}. An $F$ with~\Hyp\ that takes a negative value has $p>1$, by
\cref{prop:small} and \cref{cor:low}\,(b), and it is polyconvex if $p=2$, by
\cref{thm:iff}\,(b). \Cref{sec:signed} treats the remaining exponents $p>1$, $p\ne2$.

\subsection{Related work}\label{sec:intro-related}

Guerra and Kristensen~\cite{GuerraKristensen2022} recall two results related to their question.
Rank-one convexity implies polyconvexity for compact connected isotropic subsets of $\M$
\cite{CardaliaguetTahraoui2002a,CardaliaguetTahraoui2002b,ContiDeLellisMuellerRomeo2003}, and for
isotropic $0$-homogeneous integrands on matrices of positive determinant
\cite{MartinGhibaNeff2017}. The reduction of the polyconvexity of an isotropic function on $\M$
to supports on diagonal matrices goes back to Dacorogna and
Koshigoe~\cite[Theorem~2.1]{DacorognaKoshigoe1993}; \cref{lem:lift} is a form of it. We are not
aware of an earlier proof of \cref{thm:main,thm:closed,thm:iff,thm:low}.

The Dacorogna--Marcellini integrands
$f_\gamma(A)=\norm A^2(\norm A^2-2\gamma\det A)$, with $\norm A$ the Euclidean norm, satisfy
\ref{H1} and \ref{H2} with $p=4$. They are rank-one convex if and only if
$\abs\gamma\le2/\sqrt3$ and polyconvex if and only if $\abs\gamma\le1$
\cite{DacorognaMarcellini1988,AlibertDacorogna1992}. Since $\norm A^2\ge2\abs{\det A}$, with
equality on $\CO$, $f_\gamma$ is nonnegative exactly when $\abs\gamma\le1$.
\Cref{thm:iff}\,(a) explains the second threshold: for $p>2$, polyconvexity is rank-one
convexity together with nonnegativity. By \cref{thm:A}, $f_\gamma$ is quasiconvex for
$\abs\gamma\le2/\sqrt3$. The integrands $f_\gamma$ with $1<\abs\gamma\le2/\sqrt3$ are the
subject of the preprint~\cite{BrunoPasqualotto2026}.

Cassese~\cite[Conjecture~1]{Cassese2026} conjectures that an integrand with~\Hyp\ that changes
sign is quasiconvex at $0$; the conjecture is stated for $f(tA)=\abs t^pf(A)$. This condition
implies~\ref{H2}, so for integrands on $\M$ \cref{thm:A} gives the conclusion of the conjecture,
and quasiconvexity at every matrix. The Burkholder function is an integrand of this kind.

\subsection{The strategy}\label{sec:strategy}

Identify $\R^2$ with $\C$. A matrix acts as $\zeta\mapsto z\zeta+w\bar\zeta$ for two complex
numbers $z$, $w$, its conformal and its anticonformal part. In these coordinates
$\det A=\abs z^2-\abs w^2$ and $\abs A=\abs z+\abs w$, a rotation on either side multiplies $z$
and $w$ by numbers of modulus one, and a line is a rank-one line exactly when $z$ and $w$ move
with the same speed. So an isotropic $F$ is a function $\Phi(\abs z,\abs w)$ of two nonnegative
variables, its \emph{profile}, and rank-one convexity of $F$ becomes convexity of $\Phi$ along
the images of those lines. The proofs use only this and the homogeneity of $\Phi$.

On the segment $\abs z+\abs w=1$ the profile is the convex function $h$. A line below $h$ is the
same as a function $\abs A^{p-2}(c_0\abs A^2+m\det A)$ below $F$: these \emph{Burkholder-type
functions} are the ones whose restriction to the segment is affine. So every supporting line of
$h$ gives a Burkholder-type minorant of $F$ that touches $F$ along a ray. Where such a minorant
touches, rank-one convexity gives one inequality of first order along every rank-one line
through the point of contact (\cref{prop:tangent}). The rest of each proof is algebra with this
inequality.
\begin{enumerate}[label=\textup{(\arabic*)}]
\item \emph{$p\ge2$, $F\ge0$.} Let $\kappa=(p-2)/p$. If the supporting line of $h$ at $t$ is
  nonnegative at $\pm\kappa$, the Burkholder-type minorant has a polyaffine minorant that touches
  at the same point (\cref{lem:burkholder}); this gives the support with $c^\ast=(p-1)m$.
  Otherwise the candidate support $q$ is negative at the origin. By homogeneity it is then enough
  to prove $q\le\Phi$ at the points where $\lambda\mapsto\lambda^pq(\cdot/\lambda)$ is
  stationary. In the coordinates $(\abs w,\det A)$ these points form a line. This line is the
  image of a rank-one line through the base point, and on it $q\le\Phi$ is the first-order
  inequality.
\item \emph{$1\le p<2$, $F\ge0$.} The supporting lines of $h$ at the ends $t=\pm1$ are multiples
  of the functions $h$ of $B_p$ and of its mirror image. Every other supporting line is squeezed
  between them, which is the condition $(2-p)\abs m\le pc_0$.
\item \emph{$0<p<1$.} Along a rank-one line $F$ is convex and grows at most like the power $p$
  of the distance, hence is constant.
\item \emph{$p>1$, $p\ne2$, $F$ changes sign.} After a reflection, $h(-1)<0$, the negative values
  of $F$ are at matrices of negative determinant, and all supporting slopes of $h$ are positive.
  Let $a=\abs{p-2}/p$. Where the supporting line $\ell$ of $h$ satisfies $\ell(-a)\ge0$, the
  Burkholder-type minorant is $mB_p+\ell(-a)\abs A^p$, which is quasiconvex. At the other points,
  the \emph{hard bases}, the argument of item~(1) gives a function $\alpha+v\abs w+c\det A$ (for
  $p>2$) or $\alpha+v\abs z+c\det A$ (for $p<2$) that touches $F$ and lies below it on the unit
  ball and on a cone.

  To go further we integrate scaled truncated determinants against a convex weight. Let $N$ be
  continuous and convex on $[0,\infty)$ with $N(0)\le0$, and let $e=\pm1$. With
  $\sigma=\abs z+\abs w=\abs X$ and $\delta=\abs z-\abs w$, the function
  \[
    \Psi^e_N(X)=(\sigma+e\delta)\,N(\sigma)-2\int_0^\sigma N(r)\,dr
  \]
  is quasiconvex (\cref{prop:radial}). For the hinge function $N(r)=(r-\rho)_+$ it is
  $\rho^2K_e(X/\rho)$, with $K_1(X)=L(JX)+\det X$ and $K_{-1}(X)=L(X)-\det X$, and a convex $N$
  is a limit of sums of an affine function and of nonnegative combinations of hinge functions.
  With $N(r)=\pm r_+^{p-1}$, $\Psi^e_N$ is $B_p$ or its mirror image (\cref{cor:burkholder}), as
  in the integration over scales of Baernstein and
  Montgomery-Smith~\cite{BaernsteinMontgomerySmith1999}; this gives the quasiconvexity of $B_p$
  used above. The new step is the choice of
  $N$ adapted to $F$ at each hard base. Following the supporting lines of $h$ from the base,
  with a scale $S$ given by a quadrature, produces a convex $N$ such that $\Psi^e_N$ plus a
  constant lies below $F$ on a ball of any given radius and touches $F$ at the base
  (\cref{prop:minorant,cor:minorant}).
\end{enumerate}

\subsection{Formal verification}\label{sec:intro-formal}

The results of this paper are formalized in the proof assistant Lean~4~\cite{Lean4}, using its
mathematical library Mathlib~\cite{Mathlib20}. The proofs are in the repository~\cite{Companion},
and Lean's kernel checks every step. \Cref{thm:A}, \cref{thm:main,thm:closed,thm:iff,thm:low},
\cref{cor:equiv,cor:low} and \cref{prop:small,prop:radial} are stated in a file that imports
only Mathlib, and their proofs use only Lean's standard axioms. The formal proofs of
\cref{thm:A}, \cref{cor:low}\,(a) and \cref{prop:radial} take the quasiconvexity of the truncated
determinant from the formalization~\cite{IwaniecLean} of~\cite{SharpHI}. \Cref{app:formal}
describes the formal statements and lists the corresponding Lean declarations.

\subsection{Organization}\label{sec:intro-organization}

\Cref{sec:setting} introduces the conformal coordinates, the profile and its supports.
\Cref{sec:roc} states what rank-one convexity gives for the profile. \Cref{sec:pge2} constructs
the supports for $p\ge2$, \cref{sec:plt2} treats the exponents below $2$, and \cref{sec:signed}
the integrands that change sign; it completes the proof of \cref{thm:A}. \Cref{sec:discussion}
collects remarks and open questions. These sections give the statements and the ideas of the
proofs. The appendices supply the details in the same order:
\cref{app:setting,app:roc,app:pge2,app:plt2,app:signed} prove the statements of
\cref{sec:setting,sec:roc,sec:pge2,sec:plt2,sec:signed}, and \cref{app:formal} describes the
formal verification.

\section{Setting: coordinates, profiles and supports}\label{sec:setting}

This section turns the problem for $F$ on $\M$ into a problem for a function of two nonnegative
variables. The proofs are in \cref{app:setting}.

\subsection{Conformal coordinates}\label{sec:coordinates}

Identify $\R^2$ with $\C$. Every $A=(a_{ij})\in\M$ acts as $A\zeta=z\zeta+w\bar\zeta$ with
\[
  z=\tfrac12\bigl((a_{11}+a_{22})+i(a_{21}-a_{12})\bigr),\qquad
  w=\tfrac12\bigl((a_{11}-a_{22})+i(a_{21}+a_{12})\bigr),
\]
and $A\mapsto(z,w)$ is an $\R$-linear isomorphism $\M\to\C^2$. We call $z$ the \emph{conformal}
and $w$ the \emph{anticonformal part} of $A$, and $(z,w)$ its \emph{coordinates}.

\begin{lemma}\label{lem:coordinates}
\begin{enumerate}[label=\textup{(\alph*)}]
\item $\det A=\abs z^2-\abs w^2$ and $\abs A=\abs z+\abs w$.
\item A matrix $B\ne0$ with coordinates $(\alpha,\beta)$ has rank one if and only if
  $\abs\alpha=\abs\beta$.
\item If $R_\vartheta$ denotes the rotation by the angle $\vartheta$, the matrix
  $R_\vartheta AR_{\vartheta'}$ has the coordinates
  $(e^{i(\vartheta+\vartheta')}z,e^{i(\vartheta-\vartheta')}w)$.
\item $JA$, with $J=\diag(1,-1)$, has the coordinates $(\bar w,\bar z)$.
\end{enumerate}
\end{lemma}

By (c), the two rotations act on the phases of $z$ and $w$ independently. So $F$ is isotropic if
and only if $F(A)=\Phi(\abs z,\abs w)$, where
\[
  \Phi(x,y)=F(\diag(x+y,x-y)),\qquad x,y\ge0,
\]
is the \emph{profile} of $F$; the matrix $\diag(x+y,x-y)$ has the coordinates $(x,y)$. In the
quadrant of the variables $(x,y)=(\abs z,\abs w)$, the axis $y=0$ carries the conformal
matrices, the axis $x=0$ the anticonformal ones, and the diagonal $x=y$ the matrices of rank at
most one. The two axes together are the set $\CO$ of the Introduction.

\subsection{Profiles}

\begin{lemma}\label{lem:profile}
Let $F$ satisfy~\Hyp, with profile $\Phi$. Then
\begin{enumerate}[label=\textup{(P\arabic*)},leftmargin=3em]
\item\label{P1} $\Phi(\lambda x,\lambda y)=\lambda^p\Phi(x,y)$ for $\lambda>0$ and $x,y\ge0$;
\item\label{P2} for all $z_0,w_0,\alpha,\beta\in\C$ with $\abs\alpha=\abs\beta$, the function
  $\tau\mapsto\Phi(\abs{z_0+\tau\alpha},\abs{w_0+\tau\beta})$ is convex on $\R$.
\end{enumerate}
\end{lemma}

Property \ref{P2} is the rank-one convexity of $F$, by \cref{lem:coordinates}\,(b). From now on
a \emph{profile of degree $p$} is a function $\Phi$ on $[0,\infty)^2$ with \ref{P1} and
\ref{P2}; all the arguments use only these two properties. If $\Phi$ is a profile, so is its
\emph{mirror image} $\Phi^\ast(x,y)=\Phi(y,x)$; by \cref{lem:coordinates}\,(d) it is the profile
of $A\mapsto F(JA)$. For $p>0$, \ref{P1} gives $\Phi(0,0)=0$.

By \ref{P1} a profile is determined by its values on the segment $x+y=1$, which we parametrize
by
\[
  (x_t,y_t)=\Bigl(\frac{1+t}2,\frac{1-t}2\Bigr),\qquad -1\le t\le1 .
\]
The ends $t=1$ and $t=-1$ are the conformal and the anticonformal point, and $t=0$ is the point
of rank one. Let
\[
  h(s)=\Phi(x_s,y_s),\qquad -1\le s\le1 .
\]
With $\sigma=x+y$ and $\delta=x-y$, \ref{P1} gives $\Phi(x,y)=\sigma^ph(\delta/\sigma)$ for
$\sigma>0$. For the profile of $F$ this reads $h(s)=F(\diag(1,s))$ and
$F(A)=\abs A^p\,h(\det A/\abs A^2)$, since $\sigma=\abs A$ and $\sigma\delta=\det A$.

\subsection{Supports}

A \emph{polyaffine support} of $F$ at $A_0$ is a polyaffine function $P$ with $P\le F$ and
$P(A_0)=F(A_0)$.

\begin{lemma}\label{lem:supports}
Let $F\colon\M\to\R$.
\begin{enumerate}[label=\textup{(\alph*)}]
\item If $F$ has a polyaffine support at $A_0$, then
  $F(A_0)\le\int_{(0,1)^2}F(A_0+\nabla\varphi)\,dx$ for every Lipschitz $\varphi$ vanishing
  outside $(0,1)^2$ for which the integral exists.
\item $F$ is polyconvex if and only if it has a polyaffine support at every matrix.
\end{enumerate}
In particular polyconvex functions are quasiconvex, continuous and rank-one convex.
\end{lemma}

Part (a) holds because polyaffine functions are null Lagrangians. More generally, if $G\le F$
with $G(A_0)=F(A_0)$ and $G$ is quasiconvex at $A_0$, that is, satisfies the inequality of
quasiconvexity at $A_0$ for every test map, then $F$ is quasiconvex at $A_0$. The quasiconvexity
statements of \cref{sec:pge2,sec:plt2} are proved in this way, with a minorant that touches. In
\cref{sec:signed} the minorant lies below $F$ on a bounded set of matrices that contains the
values of $A_0+\nabla\varphi$, which is enough for a given test map $\varphi$.

\begin{lemma}\label{lem:lift}
Let $F$ be isotropic with profile $\Phi$, and let $A_0$ have the coordinates $(z_0,w_0)$.
Suppose that $k,b_1,b_2,c\in\R$ with $b_1,b_2\ge0$ satisfy
\begin{equation}\label{eq:profilesupport}
  k+b_1x+b_2y+c(x^2-y^2)\le\Phi(x,y)\qquad\text{for all }x,y\ge0,
\end{equation}
with equality at $(\abs{z_0},\abs{w_0})$. Then $F$ has a polyaffine support at $A_0$.
\end{lemma}

We call a function as in \eqref{eq:profilesupport}, with $b_1,b_2\ge0$, a \emph{support of the
profile}. The sign condition on $b_1$ and $b_2$ lets the support pass from the moduli to the
matrix, because $b_1\abs z$ dominates the linear function $b_1\re(uz)$ for every $u$ of modulus
one.
Supports scale: if \eqref{eq:profilesupport} holds with equality at $(x_0,y_0)$, then for
$\sigma>0$ the function with the coefficients
$(\sigma^pk,\sigma^{p-1}b_1,\sigma^{p-1}b_2,\sigma^{p-2}c)$ is a support with equality at
$(\sigma x_0,\sigma y_0)$, by \ref{P1}. So it is enough to find supports of the profile at the
points $(x_t,y_t)$ and at the origin.

\section{What rank-one convexity gives}\label{sec:roc}

Throughout this section $\Phi$ is a profile of degree $p>0$. The section collects what the
constructions use from \ref{P2}. The proofs are in \cref{app:roc}.

\subsection{Two special lines}

\begin{lemma}\label{lem:basic}
\begin{enumerate}[label=\textup{(\alph*)}]
\item The function $\tilde h(s)=\Phi\bigl(\abs{1+s}/2,\abs{1-s}/2\bigr)$ is convex on $\R$;
  it equals $h$ on $[-1,1]$.
\item If $0\le x_0\le x_1$, $0\le y_0\le y_1$ and $x_1^2-x_0^2=y_1^2-y_0^2$, then
  $\Phi(x_0,y_0)\le\Phi(x_1,y_1)$.
\end{enumerate}
\end{lemma}

Part (a) is convexity along the line $s\mapsto\diag(1,s)$. Part (b) says that $\Phi$ grows along
the hyperbolas $x^2-y^2=\mathrm{const}$, on which the determinant is constant; it is the
Baker--Ericksen inequality in integrated form, and it comes from the line
$\tau\mapsto(x_0+i\tau,y_0+i\tau)$, along which the profile is convex and even.

\subsection{Burkholder-type minorants}

\begin{definition}\label{def:minorant}
For $c_0,m\in\R$ the \emph{Burkholder-type function} with the coefficients $(c_0,m)$ is
\[
  G_{c_0,m}(x,y)=(x+y)^{p-1}\bigl(c_0(x+y)+m(x-y)\bigr),\qquad x,y\ge0 .
\]
\end{definition}

It is the profile of $A\mapsto\abs A^{p-2}(c_0\abs A^2+m\det A)$, and
$G_{c_0,m}(x,y)=\sigma^p\ell(\delta/\sigma)$ with the affine function $\ell(s)=c_0+ms$: the
Burkholder-type functions are the homogeneous functions that are affine on the segment
$x+y=1$. The Burkholder function is $B_p=G_{\abs\kappa,1}$ with $\kappa=(p-2)/p$. Hence
$G_{c_0,m}\le\Phi$ on the quadrant if and only if $\ell\le h$ on $[-1,1]$. Since $\tilde h$ is
convex on $\R$, $h$ has a supporting line at every $t\in[-1,1]$:

\begin{lemma}\label{lem:minorant}
For every $t\in[-1,1]$ there is $m\in\R$ such that, with $H=h(t)$, the Burkholder-type function
$G_{H-mt,m}$ lies below $\Phi$ on the quadrant and equals $\Phi$ at $(x_t,y_t)$.
\end{lemma}

\subsection{The first-order inequality}

A convex function lies above its tangent lines. For a profile, which need not be differentiable,
the minorant supplies the tangent.

\begin{proposition}[The first-order inequality]\label{prop:tangent}
Let $\abs t<1$, $H,m\in\R$, and suppose that $G_{H-mt,m}\le\Phi$ on the quadrant and
$\Phi(x_t,y_t)=H$. Put
\[
  \Phi_x=pH+(1-t)m,\qquad\Phi_y=pH-(1+t)m .
\]
Then for all $\alpha,\beta\in\C$ with $\abs\alpha=\abs\beta$ and all $\tau\in\R$,
\begin{equation}\label{eq:tangent}
  \Phi\bigl(\abs{x_t+\tau\alpha},\abs{y_t+\tau\beta}\bigr)\ge
  H+\tau\bigl(\Phi_x\re\alpha+\Phi_y\re\beta\bigr).
\end{equation}
\end{proposition}

The numbers $\Phi_x$ and $\Phi_y$ are the partial derivatives at $(x_t,y_t)$ of the minorant
$G_{H-mt,m}$; they are those of $\Phi$ where $\Phi$ is differentiable. The proof compares the
convex function $\tau\mapsto\Phi(\abs{x_t+\tau\alpha},\abs{y_t+\tau\beta})$ with the same
function of the minorant, which lies below it, agrees with it at $\tau=0$ and is differentiable
there; a convex function lies above the tangent of such a comparison function
(\cref{lem:convex}).

The inequality \eqref{eq:tangent} holds for every real $\tau$, and the constructions use it far
from $\tau=0$, on two lines. The first joins the base matrix to a rotated copy of it, where
$\Phi$ has the same value $H$:

\begin{corollary}[Baker--Ericksen]\label{cor:BE}
In the setting of \cref{prop:tangent}, $pH-2tm\ge0$.
\end{corollary}

The second is the line of \cref{lem:line} below.

\subsection{The supporting lines at the ends}

\begin{lemma}\label{lem:ends}
For every $s\in\R$,
\[
  \tilde h(s)\ge h(1)\Bigl(1+\frac p2(s-1)\Bigr)\qquad\text{and}\qquad
  \tilde h(s)\ge h(-1)\Bigl(1-\frac p2(s+1)\Bigr).
\]
\end{lemma}

At the conformal point the slope of $h$ is therefore at most $\tfrac p2h(1)$. For $p\le2$ the
first inequality says that $F\ge F(I)\,B_p/B_p(I)$ when $\Phi$ is the profile of $F$: the line
$s\mapsto1+\tfrac p2(s-1)$ is, up to the factor $B_p(I)=2/p$, the function $h$ of $B_p$. This is
the extremality of the Burkholder function; cf.~\cite[Theorem~4.2]{GuerraKristensen2022}
and~\cite{AIPS2015}.

\section{Supports for \texorpdfstring{$p\ge2$}{p >= 2}}\label{sec:pge2}

In this section $p\ge2$ and $\Phi$ is a profile of degree $p$ with $\Phi\ge0$ on the quadrant.
We construct a support of the profile at every point $(x_t,y_t)$ with $\abs t<1$, then at the
ends, and deduce \cref{thm:main,thm:closed,thm:iff} and \cref{cor:equiv}. The proofs are in
\cref{app:pge2}.

Fix $\abs t<1$, write $H=h(t)$, and fix $m$ as in \cref{lem:minorant}: with $c_0=H-mt$, the
function $G_{c_0,m}$ lies below $\Phi$ and equals $\Phi$ at $(x_t,y_t)$. Let
\[
  C=pH-2tm\ge0
\]
(\cref{cor:BE}) and let $\Phi_x$, $\Phi_y$ be as in \cref{prop:tangent}. Let $\kappa=(p-2)/p$,
and let $\ell(s)=c_0+ms$ be the supporting line of $h$ at $t$, so that
$p\,\ell(\mp\kappa)=pc_0\mp(p-2)m$. The number $\kappa$ is where the function $h$ of the
Burkholder function $B_p$ vanishes: $B_p$ is negative exactly where $\det A<-\kappa\abs A^2$.
There are three cases, according to the sign of the supporting line at $\pm\kappa$:
\begin{align*}
  &\text{III:}\quad (p-2)\abs m\le pc_0,&&\text{that is, }\ell(\kappa)\ge0\text{ and }\ell(-\kappa)\ge0;\\
  &\text{I:}\quad pc_0-(p-2)m<0,&&\text{that is, }\ell(-\kappa)<0;\\
  &\text{II:}\quad pc_0+(p-2)m<0,&&\text{that is, }\ell(\kappa)<0.
\end{align*}
In Case III the minorant $G_{c_0,m}$ is a nonnegative combination of $B_p$, of its mirror image
and of $\abs A^p$, and a polyaffine function below its positive part touches it at the same
point (\cref{lem:burkholder}). In Cases I and II the support comes from the first-order
inequality along one rank-one line.

\subsection{Case III: an inequality of Burkholder type}\label{sec:caseIII}

\begin{lemma}\label{lem:burkholder}
Let $p\ge2$, and let $c_0,m\in\R$ satisfy $(p-2)\abs m\le pc_0$. Then for all $\sigma\ge0$ and
$\delta$ with $\abs\delta\le\sigma$,
\begin{equation}\label{eq:burkholder}
  -(p-1)c_0+pc_0\sigma-(p-2)m\delta+(p-1)m\sigma\delta\le
  \max\bigl\{0,\sigma^{p-1}(c_0\sigma+m\delta)\bigr\}.
\end{equation}
\end{lemma}

With $\sigma=x+y$ and $\delta=x-y$ the right side is the positive part of $G_{c_0,m}$, and the
left side is a support of the profile in the sense of \eqref{eq:profilesupport}. For
$(c_0,m)=(\kappa,1)$ the right side is the positive part of the Burkholder function $B_p$, whose
polyconvexity for $p\ge2$ is \cite[Proposition~1.6]{GuerraKristensen2022}. The proof of the
lemma uses two polynomial identities and Bernoulli's inequality
$\sigma^p\ge\tfrac p2\sigma^2-(\tfrac p2-1)$.

In Case III, since $\Phi\ge0$ and $\Phi\ge G_{c_0,m}$, the lemma gives, for $x,y\ge0$,
\begin{equation}\label{eq:caseIII}
  -(p-1)c_0+(pc_0-(p-2)m)x+(pc_0+(p-2)m)y+(p-1)m(x^2-y^2)\le\Phi(x,y).
\end{equation}
At $(x_t,y_t)$ the left side is $c_0+mt=H$. So \eqref{eq:caseIII} is a support of the profile at
$(x_t,y_t)$, with the determinant coefficient $(p-1)m$. For $p=2$ we are always in Case III,
because $c_0=\tfrac12C\ge0$, and \eqref{eq:caseIII} holds without the assumption $\Phi\ge0$:
it reduces to $c_0(\sigma-1)^2\ge0$.

\subsection{Case I: one rank-one line}\label{sec:caseI}

Let $p>2$, $H>0$ and
\begin{equation}\label{eq:caseIhyp}
  p(H-mt)-(p-2)m<0 .
\end{equation}
Then $m>0$ (\cref{lem:caseI-algebra}). Define
\[
  c=\frac{\Phi_x}{1+t}>0,\qquad\gamma=\frac{2C}{1+t}\ge0,\qquad k=ct-(p-1)H<0 .
\]

\begin{proposition}\label{prop:caseI}
In Case I, $k+\gamma y+c(x^2-y^2)\le\Phi(x,y)$ for all $x,y\ge0$, with equality at
$(x_t,y_t)$.
\end{proposition}

The candidate $q(x,y)=k+\gamma y+c(x^2-y^2)$ is the function of the form
\eqref{eq:profilesupport} with $b_1=0$ that is tangent to the minorant at $(x_t,y_t)$. Here
\cref{lem:burkholder} does not apply, and the proof uses the homogeneity of $\Phi$ instead. For
a fixed point $(x,y)$ consider
\[
  \psi(l)=l^p\,q(x/l,y/l)=kl^p+\gamma yl^{p-1}+c(x^2-y^2)l^{p-2},\qquad l>0 .
\]
By \ref{P1}, $q\le\Phi$ at $(x/l,y/l)$ means $\psi(l)\le\Phi(x,y)$, and the claim is
$\psi(1)\le\Phi(x,y)$. Since $k<0$ and $p>2$, $\psi$ tends to $0$ at $0$ and to $-\infty$ at
infinity. If $\psi(1)\le0$ there is nothing to prove, because $\Phi\ge0$. Otherwise $\psi$ has a
positive maximum, and it is enough to prove $q\le\Phi$ at the point $(x/l,y/l)$ where the maximum
is attained:

\begin{lemma}\label{lem:psi}
Let $p>2$, $a_0<0$ and $a_1,a_2\in\R$, and let $\psi(l)=a_0l^p+a_1l^{p-1}+a_2l^{p-2}$ for
$l>0$. If $\psi(1)>0$, there is $l>0$ with $pa_0l^2+(p-1)a_1l+(p-2)a_2=0$ and
$\psi(l)\ge\psi(1)$.
\end{lemma}

So only the points $(X,Y)$ that satisfy the stationarity relation
\begin{equation}\label{eq:stationary}
  pk+(p-1)\gamma Y+(p-2)c(X^2-Y^2)=0
\end{equation}
matter. In the coordinates $(Y,D)=(\abs w,\det A)$ this relation is linear: the stationary
points form a line, which passes through the base point $(y_t,t)$ and has the slope
$\mu=-(p-1)\gamma/((p-2)c)$. On the other hand, a rank-one line through the base point along
which $w$ moves radially, $\tau\mapsto(x_t+\tau e^{i\theta},y_t+\tau)$, is also a line in these
coordinates: along it $\abs w=y_t+\tau$ and $\det A=t+(2x_t\cos\theta-2y_t)\tau$. Its slope
$2x_t\cos\theta-2y_t$ takes every value in $[-2,2t]$. By \cref{lem:mu}, $\mu$ lies in this
interval: the lower bound $\mu\ge-2$ is the hypothesis \eqref{eq:caseIhyp}, and the upper bound
follows from the Baker--Ericksen inequality and $m>0$. The stationary line is therefore the image
of a rank-one line, and on it $q\le\Phi$ is the first-order inequality:

\begin{lemma}[The stationary line]\label{lem:line}
In Case I, let $X,Y\ge0$ satisfy \eqref{eq:stationary}. Then
$k+\gamma Y+c(X^2-Y^2)\le\Phi(X,Y)$.
\end{lemma}

\Cref{prop:caseI} follows: with $l$ from \cref{lem:psi} and $(X,Y)=(x/l,y/l)$,
\[
  \Phi(x,y)=l^p\Phi(X,Y)\ge l^p\bigl(k+\gamma Y+c(X^2-Y^2)\bigr)=\psi(l)\ge\psi(1).
\]

\subsection{Case II and the three cases together}\label{sec:clamp}

Case II is the mirror image of Case I: the mirror image $\Phi^\ast$ is a nonnegative profile,
$G_{c_0,-m}$ touches it from below at $(x_{-t},y_{-t})$, and it is in Case I. Exchanging $x$ and
$y$ in \cref{prop:caseI} gives a support with $b_2=0$. One formula gives the three supports.

\begin{proposition}\label{prop:clamp}
Let $p\ge2$, $\abs t<1$ and $H>0$, and let $L$, $U$ and $c^\ast$ be as in \cref{thm:closed}.
Then
\[
  b_1=\Phi_x-(1+t)c^\ast\ge0,\qquad b_2=\Phi_y+(1-t)c^\ast\ge0,
\]
and $c^\ast t-(p-1)H+b_1x+b_2y+c^\ast(x^2-y^2)\le\Phi(x,y)$ for all $x,y\ge0$, with equality at
$(x_t,y_t)$.
\end{proposition}

The interval $[L,U]$ is the set of determinant coefficients $c$ for which the tangent function
at $(x_t,y_t)$ has $b_1,b_2\ge0$; it is not empty because $U-L=2C/(1-t^2)\ge0$. Case III is
$L\le(p-1)m\le U$, where $c^\ast=(p-1)m$. In Case I, $(p-1)m>U$ and $c^\ast=U$, so that
$b_1=0$. In Case II, $(p-1)m<L$ and $c^\ast=L$, so that $b_2=0$.

\subsection{The ends, and the proofs of the theorems}\label{sec:proofs-pge2}

At the conformal point the hyperbolas of \cref{lem:basic}\,(b) give
$\Phi(x,y)\ge h(1)(x^2-y^2)^{p/2}$ for $x\ge y$, and Bernoulli's inequality gives the support
\begin{equation}\label{eq:end-support}
  \Phi(x,y)\ge h(1)\Bigl(1+\frac p2(x^2-y^2-1)\Bigr)\qquad\text{for all }x,y\ge0,
\end{equation}
with equality at $(1,0)$; the anticonformal point is the mirror image. These two points are the
profiles of the matrices of $\CO$, where polyconvexity is
\cite[Theorem~1.3]{GuerraKristensen2022}.

\Cref{thm:main} follows. At $(x_t,y_t)$ a support of the profile is given by \cref{prop:clamp}
if $\abs t<1$ and $h(t)>0$, by the zero function if $h(t)=0$, and by \eqref{eq:end-support} if
$t=\pm1$; the zero function is a support at the origin. By scaling there is a support at every
point of the quadrant, and \cref{lem:lift,lem:supports} give the polyaffine supports and the
polyconvexity of $F$. \Cref{thm:closed} is \cref{prop:clamp} written for the matrix
$\diag(1,t)$.

For \cref{thm:iff}\,(b), the supports of Case III and of the ends do not need the sign
condition when $p=2$. For \cref{thm:iff}\,(a), a polyconvex, positively $p$-homogeneous $F$ with
$p>2$ is nonnegative: a polyaffine support $P(A)=\langle K,A\rangle+c\det A$ at $0$ satisfies
$\lambda\langle K,A\rangle+\lambda^2c\det A\le\lambda^pF(A)$ for $\lambda>0$, which forces
$K=0$ and then $c=0$ as $\lambda\downarrow0$; so $0=P\le F$.

\section{The exponents below \texorpdfstring{$2$}{2}}\label{sec:plt2}

For $p<2$ polyaffine supports are not available, and the Burkholder-type minorants themselves
are the supports. The proofs are in \cref{app:plt2}.

\subsection{The exponents \texorpdfstring{$1\le p<2$}{1 <= p < 2}}

\begin{proposition}\label{prop:low}
Let $1\le p<2$ and let $\Phi\ge0$ be a profile of degree $p$. For every $t\in[-1,1]$ there are
$c_0,m\in\R$ with $(2-p)\abs m\le pc_0$ such that $G_{c_0,m}\le\Phi$ on the quadrant, with
equality at $(x_t,y_t)$.
\end{proposition}

Let $a=(2-p)/p\in(0,1]$. The condition $(2-p)\abs m\le pc_0$ says that the line
$\ell(s)=c_0+ms$ is nonnegative at $\pm a$, and then $G_{c_0,m}$ is a nonnegative combination of
the two extreme functions $G_{a,1}=B_p$ and $G_{a,-1}$, the mirror image of $B_p$, whose lines
vanish at $-a$ and at $a$. By \cref{lem:ends} the supporting lines of $h$ at the two ends are
multiples of these two lines: at the conformal point $h$ is supported by
$\tfrac p2h(1)(s+a)$, and at the anticonformal point by $\tfrac p2h(-1)(a-s)$. A supporting
line at an interior point $t$ is compared with them. For $t\ge-a$ its slope is at most the slope
of the chord to the conformal end, and the supporting line at that end then gives
$\ell(-a)\ge0$; for $t\le a$ the anticonformal end gives $\ell(a)\ge0$. The two ranges overlap,
and outside the overlap the missing inequality follows from the other one and from $h(t)\ge0$.

\Cref{thm:low} is \cref{prop:low} for the profile of $F$, since the Burkholder-type functions
are positively $p$-homogeneous. \Cref{cor:low}\,(a) follows by the principle of the touching
minorant: if $B_p$ is quasiconvex, so are its mirror image and every $G_{c_0,m}$ of
\cref{thm:low}.

Nonnegativity enters \cref{prop:low} at the ends, through $h(\pm1)\ge0$, and at the interior
points outside the overlap. For $p=1$ it is automatic:

\begin{lemma}\label{lem:pone}
A profile of degree $1$ is nonnegative on the quadrant.
\end{lemma}

For $p=1$ we have $a=1$, and the functions of \cref{thm:low} are
$(c_0+m)\abs z+(c_0-m)\abs w$ with $c_0\ge\abs m$. They are convex, which gives
\cref{cor:low}\,(b).

\subsection{The exponents \texorpdfstring{$0<p<1$}{0 < p < 1}}

\begin{lemma}\label{lem:sublinear}
Let $g\colon\R\to\R$ be convex with $g(\tau)/\tau\to0$ as $\tau\to+\infty$ and as
$\tau\to-\infty$. Then $g$ is constant.
\end{lemma}

Let $F$ be positively $p$-homogeneous and rank-one convex with $0<p<1$; then $F$ is continuous.
For a matrix $A$ and a matrix $B$ of rank one, homogeneity gives
$F(A+\tau B)=\tau^pF(B+A/\tau)$ for $\tau>0$, so that $F(A+\tau B)/\tau\to0$. By
\cref{lem:sublinear}, $F$ is constant along every rank-one line. The matrices of rank one span
$\M$, so $F$ is constant, and the constant is $F(0)=0$. This proves \cref{prop:small}.

\subsection{All exponents}

\Cref{cor:all} collects the cases: $F=0$ for $0<p<1$ (\cref{prop:small}), $F$ convex for $p=1$
(\cref{cor:low}\,(b)), the conditional statement for $1<p<2$ (\cref{cor:low}\,(a)), and
polyconvexity for $p\ge2$ (\cref{thm:main}).

\section{Integrands that change sign}\label{sec:signed}

This section proves \cref{thm:A} for the integrands that take a negative value. As explained
after \cref{prop:small}, such an integrand has $p>1$, and the case $p=2$ is \cref{thm:iff}\,(b).
Throughout this section $p>1$, $p\ne2$, and
\[
  e=-1\ \text{if }p<2,\qquad e=1\ \text{if }p>2,\qquad a=\frac{\abs{p-2}}p=\frac{e(p-2)}p,
\]
so that $0<a<1$; for $p<2$, $a$ is the number of \cref{sec:plt2}, and for $p>2$ it is the
number $\kappa$ of \cref{sec:pge2}. $\Phi$ is a profile of degree $p$, of any sign, and $h$,
$(x_t,y_t)$, $\sigma=x+y$ and $\delta=x-y$ are as in \cref{sec:setting}. The proofs are in
\cref{app:signed}.

\subsection{Orientation and hard bases}\label{sec:orientation}

\begin{lemma}\label{lem:orient}
Let $\Phi$ be a profile of degree $p>0$.
\begin{enumerate}[label=\textup{(\alph*)}]
\item $h(0)\ge0$.
\item $(-t)^{p/2}h(-1)\le h(t)$ for $-1\le t\le0$, and $t^{p/2}h(1)\le h(t)$ for $0\le t\le1$.
\item If $\Phi(x,y)<0$ for some $x,y\ge0$, then $h(-1)<0$ or $h(1)<0$.
\item Let $t\in[-1,1]$ and $m\in\R$ with $h(s)\ge h(t)+m(s-t)$ for $s\in[-1,1]$. If $t>-1$, then
  $m\ge-\tfrac p2h(-1)$; if $t<1$, then $m\le\tfrac p2h(1)$. In particular $m>0$ if $t>-1$ and
  $h(-1)<0$.
\end{enumerate}
\end{lemma}

Part (b) comes from the hyperbolas of \cref{lem:basic}\,(b), and part (d) from the supporting
lines at the ends (\cref{lem:ends}). By (a) and the convexity of $h$, the numbers $h(-1)$ and
$h(1)$ are not both negative. The mirror image $\Phi^\ast$ exchanges them, and for the profile
of $F$ it is the profile of $A\mapsto F(JA)$, which is quasiconvex if and only if $F$ is
(\cref{lem:reflect}). So for an integrand that changes sign we may assume $h(-1)<0$. Then
$h(1)\ge0$, hence $h\ge0$ on $[0,1]$ by (b), and $F$ is negative only at matrices of negative
determinant. By (d), every supporting line of $h$ at a point of $(-1,1]$ has a positive slope.

Let $\ell(s)=H+m(s-t)$ be a supporting line of $h$ at $t$, with $H=h(t)$. One of the two
supporting lines of \cref{lem:ends} at the ends vanishes at $-a$: the line at $-1$ vanishes at
$2/p-1$, which is $-a$ for $p>2$, and the line at $1$ vanishes at $1-2/p$, which is $-a$ for
$p<2$. As functions of matrices, $B_p=G_{a,1}$ and $\abs A^p=G_{1,0}$, so that
\begin{equation}\label{eq:easy}
  G_{H-mt,m}(A)=m\,B_p(A)+\ell(-a)\,\abs A^p
\end{equation}
(\cref{lem:easy}\,(a)).
We call $(t,m)$ an \emph{easy base} if $\ell(-a)\ge0$. If moreover $m\ge0$, which holds when
$h(-1)<0$ and $t>-1$, then \eqref{eq:easy} is a nonnegative combination of $B_p$ and of the
convex function $\abs A^p$. The other points need a new construction.

\begin{definition}\label{def:hard}
Let $\abs t<1$, $H=h(t)$, and let $m$ be the slope of a supporting line of $h$ at $t$, that is,
$h(s)\ge H+m(s-t)$ for $s\in[-1,1]$, or equivalently $G_{H-mt,m}\le\Phi$ on the quadrant. The
pair $(t,m)$ is a \emph{hard base} if this line is negative at $-a$, that is, if
$p(H-mt)-e(p-2)m<0$.
\end{definition}

\begin{lemma}\label{lem:hard}
Let $\Phi$ be a profile of degree $p$.
\begin{enumerate}[label=\textup{(\alph*)}]
\item $h(-a)\ge0$.
\item A hard base $(t,m)$ has $t<-a$ if $p<2$, and $t>-a$ if $p>2$.
\item Let $h(-1)<0$. A hard base $(t,m)$ has $m>0$, and $h(t)<0$ if $p<2$, $h(t)>0$ if $p>2$.
\end{enumerate}
\end{lemma}

Parts (a) and (b) compare a supporting line with the line of \cref{lem:ends} that vanishes at
$-a$. So, if $h(-1)<0$, a hard base is a point where $h$ is negative when $p<2$, and a point
where $h$ is positive but its supporting line is negative at $-a$ when $p>2$. At the easy bases
\eqref{eq:easy} gives the quasiconvex minorant, by the next lemma: if $m>0$ and $\ell(-a)\ge0$,
then $(p-2)m\le p(H-mt)$ for $p>2$, and $(2-p)\abs m\le p(H-mt)$ for $p<2$.

\begin{lemma}\label{lem:easy}
Let $c_0,m\in\R$, and let $G_{c_0,m}(A)=\abs A^{p-2}(c_0\abs A^2+m\det A)$.
\begin{enumerate}[label=\textup{(\alph*)}]
\item For every $p\ne0$, $G_{c_0,m}(A)=m\,B_p(A)+(c_0-\abs{1-2/p}\,m)\abs A^p$.
\item If $p>2$, $m\ge0$ and $(p-2)m\le pc_0$, then $G_{c_0,m}$ is quasiconvex.
\item If $1<p<2$ and $(2-p)\abs m\le pc_0$, then $G_{c_0,m}$ is quasiconvex.
\end{enumerate}
\end{lemma}

The quasiconvexity of $B_p$ that (b) and (c) use comes from \cref{prop:radial}
(\cref{cor:burkholder}).

\subsection{Restricted supports at a hard base}\label{sec:restricted}

For real numbers $H$, $m$, $t$ with $1+et\ne0$ put $C=pH-2tm$,
\begin{equation}\label{eq:restricted}
\begin{gathered}
  c=\frac{epH+(1-et)m}{1+et},\qquad v=\frac{2C}{1+et},\qquad\alpha=ct-(p-1)H,\\
  q(x,y)=\alpha+v\,r+c\,(x^2-y^2),
\end{gathered}
\end{equation}
where $r=(\sigma-e\delta)/2$, that is, $r=x$ if $p<2$ and $r=y$ if $p>2$. At a hard base we take
$H=h(t)$; then $C\ge0$ by \cref{cor:BE}. For $p>2$ the numbers $c$, $v$, $\alpha$ are the
numbers $c$, $\gamma$, $k$ of \cref{sec:caseI}, and $q$ is the function of \cref{prop:caseI}.
For $p<2$ the numbers $c$ and $v$ are given by the formulas for $c$ and $b_1$ of Case~II in
\cref{app:pge2}. In both cases $q$ is of the form
\eqref{eq:profilesupport} once $v\ge0$, and it touches $\Phi$ at the base point, but it need not
lie below $\Phi$ on the whole quadrant.

\begin{lemma}\label{lem:rs-identities}
Let $p,e,H,m,t\in\R$ with $1+et\ne0$, and let $c$, $v$, $\alpha$, $q$ be given by
\eqref{eq:restricted} with these numbers.
\begin{enumerate}[label=\textup{(\alph*)}]
\item $q(x_t,y_t)=H$, and $q(x,y)=H-tm+m(x-y)$ if $x+y=1$.
\item $c-ev/2=m$ and $\alpha+v/2=H-tm$.
\item For $x,y\ge0$ with $s=x+y>0$,
  $\alpha s^p+v\,r\,s^{p-1}+c\,(x^2-y^2)\,s^{p-2}=G_{H-mt,m}(x,y)$.
\item Let $p>1$, $p\ne2$, and $e$ as above, and let $(t,m)$ be a hard base and $H=h(t)$. Then
  $v\ge0$, and if $h(-1)<0$, then $c>0$ and $\alpha<0$.
\end{enumerate}
\end{lemma}

Part (c) says that the function $l\mapsto l^pq(x/l,y/l)$ equals, at $l=x+y$, the Burkholder-type
minorant at $(x,y)$.

\begin{proposition}[Restricted supports]\label{prop:restricted}
Let $h(-1)<0$, let $(t,m)$ be a hard base, and let $q$ be as in \eqref{eq:restricted} with
$H=h(t)$. Then $q(x,y)\le\Phi(x,y)$ at the points $(x,y)$ of the quadrant
\begin{enumerate}[label=\textup{(\alph*)}]
\item with $x<y$, if $p<2$;
\item with $\Phi(x,y)\ge0$, if $p>2$;
\item with $x+y\le1$.
\end{enumerate}
\end{proposition}

The proof follows that of \cref{prop:caseI}. For a point $(x,y)$ let
$\psi(l)=l^pq(x/l,y/l)=\alpha l^p+v\,r\,l^{p-1}+c\,(x^2-y^2)\,l^{p-2}$, so that the claim is
$\psi(1)\le\Phi(x,y)$, and $q\le\Phi$ at $(x/l,y/l)$ means $\psi(l)\le\Phi(x,y)$.
\begin{itemize}
\item For $p>2$ at a point where $\Phi\ge0$, the proof of \cref{prop:caseI} applies: it uses
  the sign of $\Phi$ only at the point $(x,y)$.
\item For $p<2$ and $x<y$, both $\alpha$ and $c(x^2-y^2)$ are negative, so $\psi$ tends to
  $-\infty$ at $0$ and at $\infty$ and has a maximum at a stationary point. As in
  \cref{lem:line}, the stationary points form a line in the coordinates $(\abs z,\det A)$, and
  this line is the image of a rank-one line through the base point, now one along which $z$
  moves radially: $\tau\mapsto(x_t+\tau,y_t+\tau e^{i\theta})$. On it $q\le\Phi$ is the
  first-order inequality.
\item On the ball $x+y=s\le1$, at the remaining points, $\psi$ is nonincreasing on $[s,1]$, so
  that $\psi(1)\le\psi(s)=G_{H-mt,m}(x,y)\le\Phi(x,y)$ by \cref{lem:rs-identities}\,(c). For
  $p<2$ this holds at every point with $\delta\ge t\sigma$, and for $p>2$ at every point with
  $\delta\le t\sigma$, which contains the points where $\Phi<0$.
\end{itemize}

The restricted supports at two hard bases, scaled to two points, are monotone in the following
sense.

\begin{proposition}[Monotonicity of the supports]\label{prop:monotone}
Let $h(-1)<0$, let $(t_1,m_1)$ and $(t_2,m_2)$ be hard bases, with the numbers $c_i$ and $v_i$
of \eqref{eq:restricted} for $H=h(t_i)$, and let $S_1,S_2>0$. Put
$\hat v_i=S_i^{p-1}v_i$, $\hat c_i=S_i^{p-2}c_i$, $r_i=S_i(1-et_i)/2$ and $d_i=S_i^2t_i$. Then
\[
  (\hat v_2-\hat v_1)(r_2-r_1)+(\hat c_2-\hat c_1)(d_2-d_1)\ge0 .
\]
\end{proposition}

The numbers $r_i$ and $d_i$ are the values of $r$ and of $x^2-y^2$ at the point
$S_i(x_{t_i},y_{t_i})$, and with the number $\alpha_i$ of \eqref{eq:restricted} the function
$S_i^p\alpha_i+\hat v_ir+\hat c_i(x^2-y^2)$ is the restricted support of the base $i$ scaled to
$S_i$ (\cref{sec:setting}). It equals $\Phi$ at
$S_i(x_{t_i},y_{t_i})$, and it lies below $\Phi$ at the base point of the other: for $p<2$ that
point has $x<y$, since $t_j<-a<0$, and for $p>2$, $\Phi>0$ there, by \cref{lem:hard}; so
\cref{prop:restricted}\,(a) or (b) applies. Adding the two inequalities gives the proposition.

\subsection{Radial integration of the truncated determinant}\label{sec:radial}

For $e\in\R$, a continuous $N\colon\R\to\R$ and a matrix $X$ with the coordinates $(z,w)$, let
$\sigma=\abs z+\abs w$, $\delta=\abs z-\abs w$, and
\[
  \Psi^e_N(X)=(\sigma+e\delta)\,N(\sigma)-2\int_0^\sigma N(r)\,dr .
\]
Here $\sigma=\abs X$, and $\sigma+e\delta$ is $2\abs z$ for $e=1$ and $2\abs w$ for $e=-1$. We
write $\Psi^e_N$ also for the same function of the profile variables,
$\Psi^e_N(x,y)=(x+y+e(x-y))\,N(x+y)-2\int_0^{x+y}N(r)\,dr$ for $x,y\ge0$, so that
$\Psi^e_N(X)=\Psi^e_N(\abs z,\abs w)$.

\begin{proposition}[Radial integration of the truncated determinant]\label{prop:radial}
Let $N\colon\R\to\R$ be a continuous function that is convex on $[0,\infty)$, with
$N(0)\le0$, and let $e\in\{1,-1\}$. Then $\Psi^e_N$ is quasiconvex on $\M$.
\end{proposition}

The building blocks are the affine functions $N$ and the hinge functions $N(r)=(r-\rho)_+$. Let
$L$ be the truncated determinant \eqref{eq:truncated}, $J=\diag(1,-1)$, and
\[
  K_e(X)=(\sigma-1)_+\,(1+e\delta).
\]

\begin{lemma}\label{lem:radial}
Let $e\in\R$, and let $X\in\M$.
\begin{enumerate}[label=\textup{(\alph*)}]
\item If $N(r)=n_0+n_1r$, then $\Psi^e_N(X)=-n_0(\sigma-e\delta)+en_1\det X$.
\item If $N(r)=(r-\rho)_+$ with $\rho>0$, then $\Psi^e_N(X)=\rho^2K_e(\rho^{-1}X)$.
\item $K_1(X)=L(JX)+\det X$ and $K_{-1}(X)=L(X)-\det X$, and $K_1$, $K_{-1}$ are quasiconvex.
\item Let $e\in\{1,-1\}$, let $N$ and $N'$ be continuous, and let $\abs{N-N'}\le\varepsilon$ on
  $[0,R]$. Then $\abs{\Psi^e_N(X)-\Psi^e_{N'}(X)}\le4R\varepsilon$ if $\sigma\le R$.
\end{enumerate}
\end{lemma}

\begin{lemma}\label{lem:hinge}
Let $R>0$, let $N$ be convex and continuous on $[0,R]$, and let $\varepsilon>0$. There are
$n_1\in\R$ and finitely many numbers $w_i\ge0$ and $\rho_i>0$ such that
\[
  \Bigl\lvert N(r)-N(0)-n_1r-\sum_iw_i\,(r-\rho_i)_+\Bigr\rvert\le\varepsilon\qquad
  \text{for }r\in[0,R].
\]
\end{lemma}

\Cref{prop:radial} follows. Let $e=\pm1$. For $N(r)=n_0+n_1r+\sum_iw_i(r-\rho_i)_+$ with $n_0\le0$,
$w_i\ge0$ and $\rho_i>0$, \cref{lem:radial} writes $\Psi^e_N$ as the sum of
$-n_0(\sigma-e\delta)+en_1\det X$, which is $-2n_0\abs w+n_1\det X$ or
$-2n_0\abs z-n_1\det X$, a convex function plus a null Lagrangian, and of the quasiconvex
functions $w_i\rho_i^2K_e(\rho_i^{-1}X)$. A general $N$ is approximated on $[0,R]$ by such
functions with the same value at $0$ (\cref{lem:hinge}), and then $\Psi^e_N$ is approximated on
the matrices with $\abs X\le R$ (\cref{lem:radial}\,(d)). For a Lipschitz test map $\varphi$ the
matrices $A+\nabla\varphi(x)$ lie in such a set, so Jensen's inequality passes to the limit.

The Burkholder function is itself a function $\Psi^e_N$. Let $N_p(r)=r_+^{p-1}$, where
$r_+=\max\{r,0\}$. Since $\int_0^\sigma r^{p-1}\,dr=\sigma^p/p$, for $p>0$,
\[
  \Psi^e_{\pm N_p}(X)=\pm\sigma^{p-1}\Bigl(\Bigl(1-\frac2p\Bigr)\sigma+e\delta\Bigr),
\]
while $B_p(X)=\sigma^{p-1}(\abs{1-2/p}\sigma+\delta)$, because $\abs X=\sigma$ and
$\det X=\sigma\delta$. The function $N_p$ is convex on $[0,\infty)$ for $p\ge2$, and $-N_p$ for
$1<p\le2$.

\begin{corollary}[The Burkholder function]\label{cor:burkholder}
Let $N_p(r)=r_+^{p-1}$.
\begin{enumerate}[label=\textup{(\alph*)}]
\item For $p\ge2$, $B_p=\Psi^1_{N_p}$ and $B_p(JX)=\Psi^{-1}_{N_p}(X)$. For $1<p\le2$,
  $B_p=\Psi^{-1}_{-N_p}$ and $B_p(JX)=\Psi^1_{-N_p}(X)$.
\item For every $p>1$, the Burkholder function $B_p$ and its mirror image $A\mapsto B_p(JA)$
  are quasiconvex.
\end{enumerate}
\end{corollary}

So the quasiconvexity of $B_p$, which is claimed as Corollary~2 of~\cite{SharpHI}, follows from
that of the truncated determinant, claimed as Theorem~1 there.

\subsection{The path of supporting lines and the radial minorant}\label{sec:path}

Let $(t_0,m_0)$ be a hard base. By \cref{prop:restricted}\,(c) its restricted support $q_0$ lies
below $\Phi$ on the unit ball. To go further we follow the supporting lines of $h$ with the
slopes $m_0-e\xi$, $\xi\ge0$, and attach to each a scale $S(\xi)$, so that the restricted
supports, scaled to $S(\xi)$, fit together into one function $\alpha_0+\Psi^e_N$.

\begin{definition}\label{def:path}
Let $t_0,m_0,R\in\R$. A \emph{path of hard bases from $(t_0,m_0)$ to the scale $R$} consists of
a number $\bar\xi\ge0$ and functions $T,k,S\colon\R\to\R$ such that, with $m(\xi)=m_0-e\xi$ and
$g(\xi)=k(\xi)-a\,m(\xi)$:
\begin{enumerate}[label=\textup{(\roman*)}]
\item for $\xi\in[0,\bar\xi]$, $(T(\xi),m(\xi))$ is a hard base and
  $k(\xi)=h(T(\xi))-m(\xi)T(\xi)$; moreover $T(0)=t_0$, and $T$ is continuous on
  $[0,\bar\xi]$;
\item at every $\xi\in[0,\bar\xi]$, $k$ has the derivative $eT(\xi)$, and $S$ has the derivative
  \[
    S'(\xi)=-\frac{S(\xi)\,(1+eT(\xi))}{p\,g(\xi)};
  \]
\item $S>0$ on $[0,\bar\xi]$, $S(0)=1$ and $S(\bar\xi)=R$.
\end{enumerate}
\end{definition}

Here $k(\xi)$ is the intercept and $g(\xi)$ the value at $-a$ of the supporting line of slope
$m(\xi)$ at $T(\xi)$. At a hard base $g<0$ and $\abs T<1$, so $S$ is increasing.

\begin{proposition}\label{prop:path}
Let $h$ be strictly convex on $[-1,1]$, and let $(t_0,m_0)$ be a hard base. For every $R\ge1$
there is a path of hard bases from $(t_0,m_0)$ to the scale $R$.
\end{proposition}

The point $T(\xi)$ is where the line of slope $m(\xi)$ supports $h$; as $\xi$ grows the slope
moves towards a supporting slope of $h$ at $-a$, where the supporting line is nonnegative at
$-a$ by \cref{lem:hard}\,(a). Up to the first $\xi$ where $g$ vanishes, every point is a hard
base, $T$ stays between $t_0$ and $-a$, and $g$ has the derivative $e(T+a)$, so that
$\abs g$ decreases at most linearly. The scale is the quadrature
$S(\xi)=\exp\int_0^\xi(1+eT)/(-pg)$, and it tends to $\infty$ where $g$ vanishes. No
orientation of $h$ is needed here.

Along a path, let $c(\xi)$, $v(\xi)$, $\alpha(\xi)$ be the numbers \eqref{eq:restricted} at the
hard base $(T(\xi),m(\xi))$, with $H=h(T(\xi))$, and put
\[
  M(\xi)=S(\xi)^{p-1}m(\xi),\qquad f(\xi)=e\,S(\xi)^{p-2}c(\xi).
\]
On the segment $x+y=S(\xi)$ the Burkholder-type function of the supporting line of slope $m(\xi)$
is $S^pk+M\delta$, and $f$ is $e$ times the determinant coefficient $\hat c$ of the restricted
support scaled to $S(\xi)$. By \cref{lem:radial}\,(a) and \cref{lem:rs-identities}\,(b), the
affine function $N_\xi(s)=eM(\xi)+f(\xi)(s-S(\xi))$ has
\[
  \Psi^e_{N_\xi}(x,y)=\hat v(\xi)\,r+\hat c(\xi)\,(x^2-y^2),\qquad
  \hat v=S^{p-1}v,\quad\hat c=S^{p-2}c,
\]
the scaled restricted support without its constant. Let
\[
  N(s)=\sup_{\xi\in[0,\bar\xi]}\bigl(eM(\xi)+f(\xi)(s-S(\xi))\bigr),\qquad
  I(s)=\alpha(0)+s\,N(s)-2\int_0^sN(r)\,dr ,
\]
the upper envelope of the lines $N_\xi$, and its function $I$.

\begin{lemma}\label{lem:path}
Let $\bar\xi$, $T$, $k$, $S$ be a path of hard bases from $(t_0,m_0)$ to the scale $R$.
\begin{enumerate}[label=\textup{(\alph*)}]
\item At every $\xi\in[0,\bar\xi]$, $M$ has the derivative $e\,f(\xi)\,S'(\xi)$.
\item If $h(-1)<0$, then $f$ is nondecreasing on $[0,\bar\xi]$.
\item If $f$ is nondecreasing on $[0,\bar\xi]$, then
  $eM(\xi')+f(\xi')(S(\xi)-S(\xi'))\le eM(\xi)$ for $\xi,\xi'\in[0,\bar\xi]$.
\end{enumerate}
\end{lemma}

The quadrature for $S$ is equivalent to (a), which makes the lines $N_\xi$ the tangents of one
convex curve: by (c), $N_{\xi'}(S(\xi))\le eM(\xi)$. Part (b) is where \cref{prop:monotone}
enters. For $\xi_1<\xi_2$ it reads, with $W=S_2+eS_1T_1>0$ and $r_i=S_i(1-eT_i)/2$,
\[
  (f_2-f_1)(S_2-S_1)\,W\ge2(r_2-r_1)\bigl(e(M_2-M_1)-f_1(S_2-S_1)\bigr),
\]
and by (a) the last factor is at most $D(S_2-S_1)$ in absolute value if $\abs{f-f_1}\le D$ on
$[\xi_1,\xi_2]$. For $\xi_2$ close to $\xi_1$, $\abs{r_2-r_1}$ is small and the inequality
improves itself: $f$ cannot decrease.

\begin{lemma}\label{lem:envelope}
Let $\bar\xi$, $T$, $k$, $S$ be a path of hard bases from $(t_0,m_0)$ to the scale $R$, and let
$N$ and $I$ be as above.
\begin{enumerate}[label=\textup{(\alph*)}]
\item $N$ is convex on $\R$.
\end{enumerate}
If $f$ is nondecreasing on $[0,\bar\xi]$, then moreover:
\begin{enumerate}[label=\textup{(\alph*)},start=2]
\item $N(S(\xi))=eM(\xi)$ for $\xi\in[0,\bar\xi]$;
\item $N(s)=ec(0)\,s-v(0)/2$ for $s\le1$;
\item $I(S(\xi))=S(\xi)^p\,k(\xi)$ for $\xi\in[0,\bar\xi]$.
\end{enumerate}
\end{lemma}

\begin{proposition}[The radial minorant]\label{prop:minorant}
Let $h(-1)<0$, and let $\bar\xi$, $T$, $k$, $S$ be a path of hard bases from $(t_0,m_0)$ to the
scale $R$. Let $\alpha_0$ be the number $\alpha$ of \eqref{eq:restricted} for $H=h(t_0)$, $m_0$
and $t_0$. There is a convex $N\colon\R\to\R$ with $N(0)\le0$ such that
\[
  \alpha_0+\Psi^e_N(x,y)\le\Phi(x,y)\qquad\text{for }x,y\ge0\text{ with }x+y\le R,
\]
with equality at $(x_{t_0},y_{t_0})$.
\end{proposition}

The function $N$ is the envelope above; since $h(-1)<0$, $f$ is nondecreasing
(\cref{lem:path}\,(b)). Put $G=\alpha_0+\Psi^e_N$. On the segment $x+y=s$, $G$ is affine in
$\delta$: $G(x,y)=I(s)+eN(s)\,\delta$. At $s=S(\xi)$, by \cref{lem:envelope}\,(b) and (d),
\[
  G(x,y)=S^pk(\xi)+S^{p-1}m(\xi)\,\delta=s^p\Bigl(k(\xi)+m(\xi)\frac\delta s\Bigr)
  \le s^p\,h\Bigl(\frac\delta s\Bigr)=\Phi(x,y).
\]
So on the sphere of radius $s\in[1,R]$, $G$ is $s^p$ times a supporting line of $h$, evaluated
at $\delta/s$. For $s\le1$, $N$ is affine by \cref{lem:envelope}\,(c), and $G=q_0$ by
\cref{lem:radial}\,(a); it lies below $\Phi$ by \cref{prop:restricted}\,(c).
\Cref{prop:path,prop:minorant} together give the next corollary.

\begin{corollary}\label{cor:minorant}
Let $h$ be strictly convex on $[-1,1]$ with $h(-1)<0$, let $(t_0,m_0)$ be a hard base, and let
$\alpha_0$ be as in \cref{prop:minorant}. For every $R\ge1$ there is a convex
$N\colon\R\to\R$ with $N(0)\le0$ such that $\alpha_0+\Psi^e_N(x,y)\le\Phi(x,y)$ for
$x,y\ge0$ with $x+y\le R$, with equality at $(x_{t_0},y_{t_0})$.
\end{corollary}

\subsection{The proof of \texorpdfstring{\cref{thm:A}}{Theorem A}}\label{sec:proofA}

\begin{lemma}\label{lem:reflect}
Each of the properties \ref{H1}, \ref{H2} and \ref{H3} passes from $F$ to $A\mapsto F(JA)$. If
$F$ is isotropic with profile $\Phi$, then $F(JA)=\Phi(\abs w,\abs z)$ for $A$ with the
coordinates $(z,w)$, and the function $h$ of the profile of $A\mapsto F(JA)$ has at $-1$ the
value $h(1)$. If $f\colon\M\to\R$ is quasiconvex, so is $A\mapsto f(JA)$.
\end{lemma}

\begin{lemma}\label{lem:perturb}
For $A$ with the coordinates $(z,w)$ let $\Theta(A)=\abs z^p+\abs w^p$.
\begin{enumerate}[label=\textup{(\alph*)}]
\item $\Theta$ is isotropic and positively $p$-homogeneous, its profile is $x^p+y^p$, and for
  $p\ge1$ it is rank-one convex.
\item Let $p>1$, let $h_0$ be convex on $[-1,1]$, and let $\varepsilon>0$. Then
  $t\mapsto h_0(t)+\varepsilon\bigl(((1+t)/2)^p+((1-t)/2)^p\bigr)$ is strictly convex on
  $[-1,1]$.
\end{enumerate}
\end{lemma}

\begin{proposition}\label{prop:interior}
Let $F$ satisfy~\Hyp\ with $p>1$, $p\ne2$, and let the function $h$ of its profile be strictly
convex on $[-1,1]$, with $h(-1)<0$. Then $F$ is quasiconvex at every matrix $A_0$ whose
coordinates $z_0$ and $w_0$ are both nonzero.
\end{proposition}

At such a matrix $t_0=(\abs{z_0}-\abs{w_0})/\abs{A_0}$ satisfies $\abs{t_0}<1$, and a supporting
slope $m_0$ of $h$ at $t_0$ is positive. At an easy base, \eqref{eq:easy} gives a quasiconvex
minorant that touches $F$ at $A_0$. At a hard base, let $\varphi$ be a test map. Its gradient is
bounded, so $\abs{A_0+\nabla\varphi}\le R\abs{A_0}$ for some $R\ge1$. With $N$ from
\cref{cor:minorant} for this $R$, the function $X\mapsto\abs{A_0}^p(\alpha_0+
\Psi^e_N(X/\abs{A_0}))$ is quasiconvex by \cref{prop:radial}, lies below $F$ at the matrices of norm
at most $R\abs{A_0}$, and equals $F$ at $A_0$. This gives Jensen's inequality for $\varphi$.

\begin{lemma}\label{lem:dense}
Let $F\colon\M\to\R$ be continuous, and quasiconvex at every matrix whose two coordinates are
nonzero. Then $F$ is quasiconvex.
\end{lemma}

\begin{proposition}\label{prop:oriented}
Let $F$ satisfy~\Hyp\ with $p>1$, $p\ne2$, and let $h(-1)<0$ for the function $h$ of its
profile. Then $F$ is quasiconvex.
\end{proposition}

For $0<\varepsilon<-h(-1)$, the integrand $F+\varepsilon\Theta$ satisfies~\Hyp, and the function
$h$ of its profile is strictly convex and negative at $-1$ (\cref{lem:perturb}). It is
quasiconvex by \cref{prop:interior,lem:dense}, and Jensen's inequality for $F$ follows as
$\varepsilon\to0$. \Cref{thm:A} follows from \cref{prop:oriented}, applied to $F$ or, by
\cref{lem:orient}\,(c) and \cref{lem:reflect}, to $A\mapsto F(JA)$, and from the cases of
\cref{sec:intro-small}.

\section{Remarks and open questions}\label{sec:discussion}

\begin{remark}[What the results rest on]
The proofs use the following inputs.
\begin{itemize}
\item \Cref{thm:main,thm:closed,thm:iff,thm:low}, \cref{cor:equiv}, \cref{cor:low}\,(b) and
  \cref{prop:small}, which contain the cases $p\le1$ and $p=2$ of \cref{thm:A}, use the
  definition of rank-one convexity, one-variable calculus, and the fact that polyaffine
  functions are null Lagrangians.
\item \Cref{thm:A}, \cref{cor:low}\,(a), \cref{prop:radial} and \cref{cor:burkholder} use in
  addition the quasiconvexity of the truncated determinant, claimed
  in~\cite[Theorem~1]{SharpHI}, which we audited with a Lean formalization~\cite{IwaniecLean}.
  \Cref{thm:A} uses it through \cref{prop:radial} at the hard bases, and through
  \cref{cor:burkholder} at the easy bases and, for nonnegative integrands with $1<p<2$, in
  \cref{cor:low}\,(a).
\end{itemize}
The quasiconvexity of the Burkholder function, which is Iwaniec's conjecture in the form
of~\cite[Conjecture~1.1]{GuerraKristensen2022} and is claimed in~\cite[Corollary~2]{SharpHI}, is
not an input: it follows from \cref{prop:radial} (\cref{cor:burkholder}). Some input of this kind is
needed, since $B_p$ satisfies~\Hyp\ and \cref{thm:A} contains its quasiconvexity.
\end{remark}

\begin{remark}[The role of the hypotheses]
Isotropy and homogeneity enter at every step, through the profile and the function $h$ of one
variable. The sign condition is not needed for quasiconvexity. It is necessary for
polyconvexity when $p>2$ (\cref{thm:iff}\,(a)), and the proof of \cref{thm:main} uses it in
three places: in Case III, where the positive part of the minorant is used; in Case I, when
$\psi(1)\le0$; and at the ends. For $1\le p<2$ it gives a Burkholder-type minorant at every
matrix (\cref{thm:low}). For an integrand that changes sign such minorants exist at the easy
bases. At a hard base $A_0$ the proof gives, for every $R\ge1$, a quasiconvex function that
touches $F$ at $A_0$ and lies below $F$ on the matrices of norm at most $R\abs{A_0}$; this is
enough because a test map has a bounded gradient.
\end{remark}

\begin{remark}[Open questions]
\Cref{thm:A} answers the question, left open by \cref{thm:iff}\,(a) for $p>2$, whether the
integrands with~\Hyp\ that take negative values are quasiconvex. The arguments do not bear on
Morrey's problem for integrands on $\M$ that are not isotropic or not homogeneous.
\end{remark}


\appendix
\crefalias{section}{appendix}
\crefalias{subsection}{appendix}
\crefname{appendix}{Appendix}{Appendices}
\Crefname{appendix}{Appendix}{Appendices}
\section{Proofs for Section~\ref{sec:setting}}\label{app:setting}

\begin{proof}[Proof of \cref{lem:coordinates}]
(a) From $a_{11}=\re z+\re w$, $a_{22}=\re z-\re w$, $a_{21}=\im z+\im w$ and
$a_{12}=\im w-\im z$ we get $\det A=\abs z^2-\abs w^2$. For $\abs\zeta=1$,
$\abs{z\zeta+w\bar\zeta}\le\abs z+\abs w$, with equality when $z\zeta$ and $w\bar\zeta$ have the
same argument; so $\abs A=\abs z+\abs w$.

(b) A nonzero matrix has rank one exactly when its determinant vanishes, that is, by (a), when
$\abs\alpha=\abs\beta$.

(c) The rotation $R_\vartheta$ is $\zeta\mapsto e^{i\vartheta}\zeta$, so
$R_\vartheta AR_{\vartheta'}\zeta=e^{i\vartheta}\bigl(ze^{i\vartheta'}\zeta
+w e^{-i\vartheta'}\bar\zeta\bigr)$.

(d) $J$ is $\zeta\mapsto\bar\zeta$, so $JA\zeta=\bar z\bar\zeta+\bar w\zeta$.
\end{proof}

\begin{proof}[Proof of \cref{lem:profile}]
\ref{P1} is the homogeneity of $F$. For \ref{P2}, the matrix with coordinates
$(z_0+\tau\alpha,w_0+\tau\beta)$ is $A_0+\tau B$, where $A_0$ and $B$ have the coordinates
$(z_0,w_0)$ and $(\alpha,\beta)$. By \cref{lem:coordinates}\,(b), $B$ has rank one, or $B=0$. In
both cases $\tau\mapsto F(A_0+\tau B)=\Phi(\abs{z_0+\tau\alpha},\abs{w_0+\tau\beta})$ is
convex.
\end{proof}

\begin{proof}[Proof of \cref{lem:supports}]
(a) Let $P$ be a polyaffine support at $A_0$. The entries and the determinant are null
Lagrangians: $\int_{(0,1)^2}P(A_0+\nabla\varphi)\,dx=P(A_0)$. Hence
$F(A_0)=P(A_0)=\int P(A_0+\nabla\varphi)\,dx\le\int F(A_0+\nabla\varphi)\,dx$.

(b) Let $T(A)=(A,\det A)$. If $F=g\circ T$ with $g$ convex, an affine support of $g$ at
$T(A_0)$, composed with $T$, is a polyaffine support of $F$ at $A_0$. Conversely, let
$\mathcal S$ be the set of affine functions $a$ on $\M\times\R$ with $a\circ T\le F$; it is
not empty, since $F$ has a polyaffine support. Let $g=\sup_{a\in\mathcal S}a$. For
$(X,d)\in\M\times\R$ let $D=\diag(d-\det X,1)$; since
$\det(X+D)+\det(X-D)=2\det X+2\det D$, the point $(X,d)$ is the midpoint of $T(X+D)$ and
$T(X-D)$, so that $a(X,d)\le\tfrac12F(X+D)+\tfrac12F(X-D)$ for $a\in\mathcal S$. Thus $g$ is
finite, and it is convex as a supremum of affine functions. Finally $g\circ T\le F$ by
definition, and $g(T(A))\ge F(A)$ by the support at $A$.

For the last statement, let $F=g\circ T$ be polyconvex. It is quasiconvex by (a) and (b). It
is continuous, because a finite convex function on $\M\times\R$ is. It is rank-one convex,
because by (b) it is a supremum of polyaffine functions, which are affine along rank-one
lines.
\end{proof}

\begin{proof}[Proof of \cref{lem:lift}]
Choose $u,v\in\C$ of modulus one with $uz_0=\abs{z_0}$ and $vw_0=\abs{w_0}$, and put
$P(A)=k+b_1\re(uz)+b_2\re(vw)+c\det A$, which is polyaffine. Since $b_1,b_2\ge0$,
$\re(uz)\le\abs z$ and $\re(vw)\le\abs w$, we get $P(A)\le\Phi(\abs z,\abs w)=F(A)$ by
\eqref{eq:profilesupport} and \cref{lem:coordinates}\,(a), with equality at $A_0$.
\end{proof}

\section{Proofs for Section~\ref{sec:roc}}\label{app:roc}

Here $\Phi$ is a profile of degree $p\in\R$; the proof of \cref{lem:minorant} uses $p>0$.

\begin{proof}[Proof of \cref{lem:basic}]
(a) is \ref{P2} for $(z_0,w_0)=(\tfrac12,\tfrac12)$ and $(\alpha,\beta)=(\tfrac12,-\tfrac12)$.
For (b), \ref{P2} with $(z_0,w_0)=(x_0,y_0)$ and $\alpha=\beta=i$ says that
$g(\tau)=\Phi\bigl(\sqrt{x_0^2+\tau^2},\sqrt{y_0^2+\tau^2}\bigr)$ is convex; it is even, so
$g(0)\le g(\tau)$. Take $\tau=\sqrt{x_1^2-x_0^2}$.
\end{proof}

\begin{proof}[Proof of \cref{lem:minorant}]
The convex function $\tilde h$ has a supporting line at $t$: there is $m$ with
$\tilde h(s)\ge H+m(s-t)$ for all $s$. On $[-1,1]$ this is $h\ge\ell$ with $\ell(s)=(H-mt)+ms$,
that is, $G_{H-mt,m}\le\Phi$ on the quadrant; at the origin both sides vanish. At $(x_t,y_t)$
the minorant has the value $\ell(t)=H$.
\end{proof}

The next lemma turns a comparison near one point into a supporting line.

\begin{lemma}\label{lem:convex}
Let $g\colon\R\to\R$ be convex and $\tau_0\in\R$. Let $k$ be a real function with $k\le g$
near $\tau_0$ and $k(\tau_0)=g(\tau_0)$, differentiable at $\tau_0$. Then
$g(\tau)\ge g(\tau_0)+k'(\tau_0)(\tau-\tau_0)$ for every $\tau\in\R$.
\end{lemma}

\begin{proof}
For $0<\varepsilon<1$, convexity gives
$g(\tau_0+\varepsilon(\tau-\tau_0))\le(1-\varepsilon)g(\tau_0)+\varepsilon g(\tau)$. For small
$\varepsilon$,
\[
  g(\tau)-g(\tau_0)\ge\frac{g(\tau_0+\varepsilon(\tau-\tau_0))-g(\tau_0)}{\varepsilon}
  \ge\frac{k(\tau_0+\varepsilon(\tau-\tau_0))-k(\tau_0)}{\varepsilon},
\]
and the right side tends to $k'(\tau_0)(\tau-\tau_0)$ as $\varepsilon\downarrow0$.
\end{proof}

\begin{proof}[Proof of \cref{prop:tangent}]
Let $X(\tau)=\abs{x_t+\tau\alpha}$ and $Y(\tau)=\abs{y_t+\tau\beta}$. The function
$g(\tau)=\Phi(X(\tau),Y(\tau))$ is convex by \ref{P2}, the function
$k(\tau)=G_{H-mt,m}(X(\tau),Y(\tau))$ satisfies $k\le g$, and $k(0)=H=g(0)$ because
$x_t+y_t=1$ and $x_t-y_t=t$. Since $x_t>0$,
$X(\tau)=\bigl((x_t+\tau\re\alpha)^2+\tau^2(\im\alpha)^2\bigr)^{1/2}$ is differentiable at $0$
with $X'(0)=\re\alpha$; likewise $Y'(0)=\re\beta$. The partial derivatives of
$G_{c_0,m}(x,y)=\sigma^{p-1}(c_0\sigma+m\delta)$ at $(x_t,y_t)$, where $\sigma=1$,
$\delta=t$ and $c_0+mt=H$, are
\[
  (p-1)H+c_0+m=\Phi_x,\qquad(p-1)H+c_0-m=\Phi_y .
\]
So $k'(0)=\Phi_x\re\alpha+\Phi_y\re\beta$, and \cref{lem:convex} gives \eqref{eq:tangent}.
\end{proof}

\begin{proof}[Proof of \cref{cor:BE}]
Let $\tau=2x_ty_t>0$, $\alpha=-y_t+i\sqrt{1-y_t^2}$ and $\beta=-x_t+i\sqrt{1-x_t^2}$, so that
$\abs\alpha=\abs\beta=1$. Then
\[
  \abs{x_t+\tau\alpha}^2=x_t^2\bigl((1-2y_t^2)^2+4y_t^2(1-y_t^2)\bigr)=x_t^2,
\]
and likewise $\abs{y_t+\tau\beta}=y_t$: the line joins the base matrix to a rotated copy of it.
By \eqref{eq:tangent},
\[
  H\ge H-2x_ty_t\,(y_t\Phi_x+x_t\Phi_y)=H-2x_ty_t\,(pH-2tm). \qedhere
\]
\end{proof}

\begin{proof}[Proof of \cref{lem:ends}]
For $u>0$, \cref{lem:basic}\,(b), from $(\sqrt u,0)$ to
$\bigl(\tfrac{1+u}2,\tfrac{\abs{1-u}}2\bigr)$, and \ref{P1} give
$\tilde h(u)\ge\Phi(\sqrt u,0)=u^{p/2}h(1)$, with equality at $u=1$. \Cref{lem:convex}
with $k(u)=u^{p/2}h(1)$ and $\tau_0=1$ gives the first inequality. The second is the first for
the mirror image.
\end{proof}

\section{Proofs for Section~\ref{sec:pge2}}\label{app:pge2}

We keep the notation of \cref{sec:pge2}: $p\ge2$, $\Phi\ge0$ is a profile of degree $p$,
$\abs t<1$, $H=h(t)$, $m$ is the slope of a supporting line of $h$ at $t$, $c_0=H-mt$,
$C=pH-2tm\ge0$, $\Phi_x=pH+(1-t)m$ and $\Phi_y=pH-(1+t)m$.

\subsection{The inequality of Burkholder type}

\begin{proof}[Proof of \cref{lem:burkholder}]
Write $E$ for the left side of \eqref{eq:burkholder} and $r=c_0\sigma+m\delta$. The hypothesis
gives $c_0\ge0$ and $(p-2)\abs{m\delta}\le pc_0\sigma$. If $\sigma=0$, then $\delta=0$ and
$E=-(p-1)c_0\le0$. Let $\sigma>0$.

\emph{Case $r\ge0$.} Bernoulli's inequality with the exponent $p/2\ge1$ gives
$\sigma^p=(\sigma^2)^{p/2}\ge\tfrac p2\sigma^2-(\tfrac p2-1)$. Together with the identity
\[
  r\Bigl(\frac p2\sigma^2-\Bigl(\frac p2-1\Bigr)\Bigr)-\sigma E
  =(\sigma-1)^2\Bigl(\frac p2c_0\sigma-\Bigl(\frac p2-1\Bigr)m\delta\Bigr)\ge0
\]
this yields $\sigma E\le r\sigma^p$, that is, $E\le\sigma^{p-1}r$.

\emph{Case $r<0$.} Let $N=-m\delta$, so that $N>c_0\sigma\ge0$. Then
\[
  4(p-1)N\sigma\,E=-\bigl(2(p-1)N\sigma-pc_0\sigma-(p-2)N\bigr)^2
  -(N-c_0\sigma)\bigl(p^2c_0\sigma-(p-2)^2N\bigr),
\]
and $(p-2)^2N\le p(p-2)N\le p^2c_0\sigma$. Hence $E\le0$.
\end{proof}

For $p=2$ the left side of \eqref{eq:caseIII} is
$-c_0+2c_0\sigma+m\sigma\delta=\sigma(c_0\sigma+m\delta)-c_0(\sigma-1)^2$, which is at most
$G_{c_0,m}=\sigma(c_0\sigma+m\delta)\le\Phi$ because $c_0=\tfrac12C\ge0$. This proves
\eqref{eq:caseIII} for $p=2$ without the assumption $\Phi\ge0$.

\subsection{Case I}

Let $p>2$ and $H>0$, and assume \eqref{eq:caseIhyp}. Recall
$c=\Phi_x/(1+t)$, $\gamma=2C/(1+t)$, $k=ct-(p-1)H$ and $\mu=-(p-1)\gamma/((p-2)c)$.

\begin{lemma}\label{lem:caseI-algebra}
In Case I, $m>0$, $c>0$, $\gamma\ge0$, and
\begin{align}
  &\gamma y_t+2ct=pH,\label{eq:I1}\\
  &c(1-t)+\Phi_y=\gamma,\label{eq:I2}\\
  &pk+(p-1)\gamma y_t+(p-2)ct=0,\label{eq:I3}\\
  &2(1+t)\,k=-\bigl((1-t)C+(p-2)(1+t)H\bigr)<0 .\label{eq:I4}
\end{align}
\end{lemma}

\begin{proof}
Subtracting $C\ge0$ from \eqref{eq:caseIhyp} gives $-m(p-2)(1+t)<0$, so $m>0$. Hence
$\Phi_x=pH+(1-t)m>0$ and $c>0$, and $\gamma\ge0$ because $C\ge0$. The four relations are
computations from the definitions; in \eqref{eq:I4} the right side is negative because $C\ge0$
and $H>0$.
\end{proof}

By \eqref{eq:I1}, $k+\gamma y_t+ct=H$: this is the equality at $(x_t,y_t)$ in
\cref{prop:caseI}.

\begin{lemma}\label{lem:mu}
In Case I, $-2\le\mu\le2t$.
\end{lemma}

\begin{proof}
From the definitions,
\begin{align*}
  (\mu+2)(p-2)c(1+t)&=-2\bigl(p(H-mt)-(p-2)m\bigr),\\
  (2t-\mu)(p-2)c(1+t)&=2\bigl(t(p-2)\Phi_x+(p-1)C\bigr).
\end{align*}
The first right side is positive by \eqref{eq:caseIhyp}. The second is nonnegative: for
$t\ge0$ both terms are, and for $t<0$ it equals
$2pH\bigl((p-1)+(p-2)t\bigr)+2(-t)m\bigl(p+(p-2)t\bigr)$, where all factors are positive.
\end{proof}

\begin{proof}[Proof of \cref{lem:line}]
Subtracting \eqref{eq:I3} from \eqref{eq:stationary} gives $X^2-Y^2=t+\mu\tau$ with
$\tau=Y-y_t$. By \cref{lem:mu} there is $\theta$ with
$\cos\theta=(\mu+1-t)/(1+t)$. Consider the rank-one line with $\alpha=e^{i\theta}$ and
$\beta=1$. Then $\abs{y_t+\tau}=Y$, and, as $2x_t=1+t$, $2y_t=1-t$ and $x_t^2-y_t^2=t$,
\[
  \abs{x_t+\tau e^{i\theta}}^2=x_t^2+\tau(\mu+1-t)+\tau^2=t+\mu\tau+(y_t+\tau)^2=X^2 .
\]
By \cref{prop:tangent}, then \eqref{eq:I2} and \eqref{eq:I1},
\[
  \Phi(X,Y)\ge H+\tau(\Phi_x\cos\theta+\Phi_y)=H+\tau(c\mu+\gamma)
  =k+\gamma Y+c(t+\mu\tau).\qedhere
\]
\end{proof}

\begin{proof}[Proof of \cref{lem:psi}]
$\psi$ extends continuously to $[0,\infty)$ with $\psi(0)=0$, and
$\psi(\Lambda)=\Lambda^{p-2}(a_0\Lambda^2+a_1\Lambda+a_2)\le0$ for
$\Lambda=\max\{1,(a_1+\abs{a_2})/\abs{a_0}\}$. So the maximum of $\psi$ on $[0,\Lambda]$ is
attained at some $l\in(0,\Lambda)$, where $\psi(l)\ge\psi(1)$ and
$0=\psi'(l)=l^{p-3}\bigl(pa_0l^2+(p-1)a_1l+(p-2)a_2\bigr)$.
\end{proof}

In the proof of \cref{prop:caseI} in \cref{sec:caseI}, \cref{lem:psi} is applied with
$a_0=k<0$, $a_1=\gamma y$ and $a_2=c(x^2-y^2)$; dividing the relation
$pkl^2+(p-1)\gamma yl+(p-2)c(x^2-y^2)=0$ by $l^2$ shows that $(X,Y)=(x/l,y/l)$ satisfies
\eqref{eq:stationary}.

\subsection{Case II and the proof of \cref{prop:clamp}}

Let $p>2$, $H>0$ and $p(H-mt)+(p-2)m<0$. The mirror image $\Phi^\ast$ is a nonnegative profile,
$G_{c_0,-m}$ touches it from below at $(x_{-t},y_{-t})$, and it is in Case I. Exchanging $x$
and $y$ in \cref{prop:caseI} gives, with
\[
  c=\frac{(1+t)m-pH}{1-t},\qquad b_1=\frac{2C}{1-t}\ge0,
\]
the support $ct-(p-1)H+b_1x+c(x^2-y^2)\le\Phi(x,y)$, with equality at $(x_t,y_t)$.

\begin{proof}[Proof of \cref{prop:clamp}]
From the definitions,
\[
  U-L=\frac{2C}{1-t^2}\ge0,\qquad U-(p-1)m=\frac{pc_0-(p-2)m}{1+t},\qquad
  (p-1)m-L=\frac{pc_0+(p-2)m}{1-t}.
\]
In Case III, $L\le(p-1)m\le U$ and $c^\ast=(p-1)m$; then $b_1=pc_0-(p-2)m$,
$b_2=pc_0+(p-2)m$ and $c^\ast t-(p-1)H=-(p-1)c_0$, and the support is \eqref{eq:caseIII}.
Otherwise $p>2$. If $m\ge0$, then $pc_0<(p-2)m$, which is Case I; there $(p-1)m>U$, so
$c^\ast=U$, $b_1=0$ and $b_2=\gamma$, and the support is that of \cref{prop:caseI}. If $m<0$,
then $pc_0<-(p-2)m$, which is Case II; there $(p-1)m<L\le U$, so $c^\ast=L$, $b_1=2C/(1-t)$ and
$b_2=0$. The equality at $(x_t,y_t)$ is a computation.
\end{proof}

\subsection{The ends}

\begin{proof}[Proof of \eqref{eq:end-support}]
For $x\ge y\ge0$, \cref{lem:basic}\,(b), from $(\sqrt{x^2-y^2},0)$ to $(x,y)$, and \ref{P1} give
$\Phi(x,y)\ge h(1)(x^2-y^2)^{p/2}$, and Bernoulli's inequality with the exponent $p/2\ge1$ gives
$(x^2-y^2)^{p/2}\ge1+\tfrac p2(x^2-y^2-1)$. For $x<y$ the right side of \eqref{eq:end-support}
is nonpositive, because $h(1)\ge0$, and $\Phi\ge0$. At $(1,0)$ both sides are $h(1)$.
\end{proof}

The mirror image gives the support $h(-1)\bigl(1-\tfrac p2(x^2-y^2+1)\bigr)$ at $(0,1)$.

\subsection{The theorems}

\begin{proof}[Proof of \cref{thm:main}]
Let $F$ be as in the theorem, with profile $\Phi$; by \cref{lem:profile}, $\Phi$ is a
nonnegative profile of degree $p$. At $(x_t,y_t)$ a support of the profile is given by
\cref{prop:clamp} if $\abs t<1$ and $h(t)>0$, by the zero function if $h(t)=0$, and by
\eqref{eq:end-support} and its mirror image if $t=\pm1$; the zero function is a support at the
origin. By scaling there is a support of the profile at every point of the quadrant, and
\cref{lem:lift} gives a polyaffine support of $F$ at every matrix. By \cref{lem:supports}, $F$
is polyconvex.
\end{proof}

\begin{proof}[Proof of \cref{thm:closed}]
The hypothesis on $m$ says that $G_{H-mt,m}$ touches $\Phi$ from below at $(x_t,y_t)$. In
coordinates, $a_{11}=\re z+\re w$ and $a_{22}=\re z-\re w$, so that
\[
  P_t(A)=c^\ast t-(p-1)H+b_1\re z+b_2\re w+c^\ast\det A
\]
with $b_1$ and $b_2$ as in \cref{prop:clamp}, and $P_t\le F$ as in the proof of \cref{lem:lift}.
At $\diag(1,t)$, where $(z,w)=(x_t,y_t)$, there is equality.
\end{proof}

\begin{proof}[Proof of \cref{thm:iff}]
(b) Let $F$ be isotropic, positively $2$-homogeneous and rank-one convex, of any sign. For
$\abs t<1$, \eqref{eq:caseIII} is a support without the sign condition. At $t=1$,
\cref{lem:ends} gives $h(s)\ge h(1)s$, so that
$h(1)(x^2-y^2)=G_{0,h(1)}(x,y)\le\Phi(x,y)$; this is a support at $(1,0)$ and at the origin. The
point $t=-1$ is the mirror image. Hence $F$ is polyconvex, as in the proof of \cref{thm:main}.
The converse is \cref{lem:supports}.

(a) By \cref{thm:main} and \cref{lem:supports} it remains to show that a polyconvex, positively
$p$-homogeneous $F$ with $p>2$ is nonnegative. Let $P(A)=\langle K,A\rangle+c\det A$ be a
polyaffine support at $0$; its constant term is $F(0)=0$. For $\lambda>0$,
\[
  \lambda\langle K,A\rangle+\lambda^2c\det A=P(\lambda A)\le\lambda^pF(A).
\]
Dividing by $\lambda$ and letting $\lambda\downarrow0$ gives $\langle K,A\rangle\le0$ for all
$A$, so $K=0$. Dividing by $\lambda^2$ and letting $\lambda\downarrow0$ gives $c\det A\le0$ for
all $A$, so $c=0$. Thus $0=P\le F$.
\end{proof}

\begin{proof}[Proof of \cref{cor:equiv}]
Rank-one convexity implies polyconvexity by \cref{thm:main}, polyconvexity implies
quasiconvexity by \cref{lem:supports}, and a Borel, locally bounded, quasiconvex function is
rank-one convex~\cite{Morrey1952,Dacorogna2008}.
\end{proof}

\section{Proofs for Section~\ref{sec:plt2}}\label{app:plt2}

\subsection{The exponents \texorpdfstring{$1\le p<2$}{1 <= p < 2}}

\begin{proof}[Proof of \cref{prop:low}]
At $t=1$, \cref{lem:ends} gives $h(s)\ge c_0+ms$ with $m=\tfrac p2h(1)\ge0$ and
$c_0=(1-\tfrac p2)h(1)$, so that $(2-p)\abs m=pc_0$; there is equality at $s=1$. The point
$t=-1$ is the mirror image.

Let $\abs t<1$, $H=h(t)\ge0$, and let $m$ be the slope of a supporting line of $\tilde h$ at
$t$: $h(s)\ge H+m(s-t)$ on $[-1,1]$. With $c_0=H-mt$ we have to show that
\[
  \mathrm{(A)}\quad pH-m\rho\ge0\qquad\text{and}\qquad\mathrm{(B)}\quad pH+m\rho'\ge0,
\]
where $\rho=pt+2-p$ and $\rho'=2-p-pt$; note that $\rho+\rho'=2(2-p)>0$. The supporting line
at $s=\pm1$ and \cref{lem:ends} at $s=t$ give
\begin{align*}
  m(1-t)&\le h(1)-H,& \rho\,h(1)&\le2H,\\
  m(1+t)&\ge H-h(-1),& \rho'h(-1)&\le2H .
\end{align*}
If $\rho\ge0$, then $(1-t)(pH-m\rho)\ge pH(1-t)-\rho(h(1)-H)=2H-\rho\,h(1)\ge0$, which is
(A). If $\rho'\ge0$, then $(1+t)(pH+m\rho')\ge pH(1+t)+\rho'(H-h(-1))=2H-\rho'h(-1)\ge0$,
which is (B). Let $\rho<0$; then $\rho'>0$ and (B) holds. If $m\ge0$, then
$pH-m\rho\ge pH\ge0$; if $m<0$, then $pH-m\rho=(pH+m\rho')-m(\rho+\rho')\ge0$. This is (A).
The case $\rho'<0$ is symmetric.
\end{proof}

In terms of $a=(2-p)/p$, $\rho\ge0$ means $t\ge-a$ and $\rho'\ge0$ means $t\le a$.

\begin{proof}[Proof of \cref{lem:pone}]
By \cref{lem:ends} with $p=1$, $h(-1)=\tilde h(-1)\ge0$ and $h(1)\ge0$, and
$h(s)\ge\tfrac12h(1)(1+s)\ge0$ on $[-1,1]$.
\end{proof}

\begin{proof}[Proof of \cref{thm:low}]
By \cref{lem:profile} the profile $\Phi$ of $F$ is a nonnegative profile of degree $p$.
The Burkholder-type functions are positively $p$-homogeneous, so
\cref{prop:low} at $t=\det A_0/\abs{A_0}^2$ gives $c_0,m$ with
$G_{c_0,m}(\abs z,\abs w)\le\Phi(\abs z,\abs w)=F(A)$ for all $A$, with equality at $A_0$; for
$A_0=0$ take $c_0=m=0$. The formula for $G_{c_0,m}$ in terms of
$B_p(A)=\tfrac2p(\abs z-(p-1)\abs w)(\abs z+\abs w)^{p-1}$ and
$B_p(JA)=\tfrac2p(\abs w-(p-1)\abs z)(\abs z+\abs w)^{p-1}$ is an identity, and its coefficients
are nonnegative because $(2-p)\abs m\le pc_0$.
\end{proof}

\begin{proof}[Proof of \cref{cor:low}]
(a) The Burkholder function $B_p$ is quasiconvex, by \cite[Corollary~2]{SharpHI} as
formalized in~\cite{IwaniecLean}; see \cref{sec:intro-low} for the test maps. So is
$A\mapsto B_p(JA)$, by composing the test maps with $J$, and both are continuous. Hence every $G_{c_0,m}$ as in \cref{thm:low} is quasiconvex,
and for a test map $\varphi$,
\[
  F(A_0)=G_{c_0,m}(A_0)\le\int_{(0,1)^2}G_{c_0,m}(A_0+\nabla\varphi)\,dx
  \le\int_{(0,1)^2}F(A_0+\nabla\varphi)\,dx .
\]

(b) By \cref{lem:pone}, $F\ge0$, so \cref{thm:low} applies with $p=1$: at every
$A_0$ there are $c_0\ge\abs m$ such that $(c_0+m)\abs z+(c_0-m)\abs w$ lies below $F$ and equals
$F$ at $A_0$. This function is convex. A function with a convex touching minorant at every
point is convex.
\end{proof}

\subsection{The exponents \texorpdfstring{$0<p<1$}{0 < p < 1}}

\begin{proof}[Proof of \cref{lem:sublinear}]
Let $s<t$. For $T>t$, convexity gives
\[
  \frac{g(t)-g(s)}{t-s}\le\frac{g(T)-g(t)}{T-t},
\]
and the right side tends to $0$ as $T\to+\infty$, since $g(T)/T\to0$. For $T<s$,
\[
  \frac{g(s)-g(T)}{s-T}\le\frac{g(t)-g(s)}{t-s},
\]
and the left side tends to $0$ as $T\to-\infty$. Hence $g(s)=g(t)$.
\end{proof}

\begin{proof}[Proof of \cref{prop:small}]
A rank-one convex function is separately convex in the entries of the matrix, hence continuous.
From \ref{H2} with $\lambda=2$, $F(0)=2^pF(0)$, so $F(0)=0$.

Let $A\in\M$ and let $B$ have rank one, and put $g(\tau)=F(A+\tau B)$, which is convex. For
$\tau>0$, homogeneity gives $g(\tau)/\tau=\tau^{p-1}F(B+A/\tau)$, which tends to $0$ as
$\tau\to+\infty$ because $p<1$ and $F$ is continuous at $B$. The same with $-B$ in place of $B$
gives $g(\tau)/\tau\to0$ as $\tau\to-\infty$. By \cref{lem:sublinear}, $g$ is constant:
$F(A+\tau B)=F(A)$ for every $\tau$.

The four matrices with one entry equal to $1$ and the others $0$ have rank one. Removing the
entries of $A$ one at a time gives $F(A)=F(0)=0$.
\end{proof}

\begin{proof}[Proof of \cref{cor:all}]
For $0<p<1$, $F=0$ by \cref{prop:small}. For $p=1$, $F$ is convex by \cref{cor:low}\,(b), hence
quasiconvex. For $1<p<2$ this is
\cref{cor:low}\,(a). For $p\ge2$, $F$ is polyconvex by \cref{thm:main}, hence quasiconvex by
\cref{lem:supports}.
\end{proof}

\section{Proofs for Section~\ref{sec:signed}}\label{app:signed}

We keep the notation of \cref{sec:signed}. The test maps are the Lipschitz maps
$\varphi\colon\R^2\to\R^2$ that vanish outside $(0,1)^2$. If $\varphi$ has the Lipschitz constant
$K$, the entries of $\nabla\varphi$ are bounded by $K$, so that the coordinates of
$\nabla\varphi(x)$ satisfy $\abs z+\abs w\le4K$: each of $\re z$, $\im z$, $\re w$, $\im w$ is
half a sum or difference of two entries. In particular $F(A+\nabla\varphi)$ is bounded and
integrable on $(0,1)^2$ if $F$ is continuous.

\subsection{Operations that preserve quasiconvexity}

\begin{lemma}\label{lem:qc-operations}
\begin{enumerate}[label=\textup{(\alph*)}]
\item If $f\colon\M\to\R$ is quasiconvex, so are $X\mapsto f(\lambda X)$ and
  $X\mapsto\kappa+f(X)$, for $\lambda,\kappa\in\R$.
\item If $g_1,\dots,g_n$ are continuous and quasiconvex and $w_1,\dots,w_n\ge0$, then
  $\sum_iw_ig_i$ is quasiconvex.
\item For $b_1,b_2\ge0$ and $c\in\R$, the function $X\mapsto b_1\abs z+b_2\abs w+c\det X$ is
  quasiconvex.
\end{enumerate}
\end{lemma}

\begin{proof}
(a) If $\varphi$ is a test map, so is $\lambda\varphi$, and
$f(\lambda(A+\nabla\varphi))=f(\lambda A+\nabla(\lambda\varphi))$. For the constant, the unit
square has measure~$1$.

(b) For a test map the functions $g_i(A+\nabla\varphi)$ are integrable, because the $g_i$ are
continuous; add the inequalities.

(c) The function is isotropic, with the profile $b_1x+b_2y+c(x^2-y^2)$, by
\cref{lem:coordinates}. This profile is a support of itself at every point, in the sense of
\eqref{eq:profilesupport}. By \cref{lem:lift} the function has a polyaffine support at every
matrix, and it is quasiconvex by \cref{lem:supports}.
\end{proof}

\subsection{Orientation and hard bases}

\begin{proof}[Proof of \cref{lem:orient}]
(a) By \ref{P2} with $z_0=w_0=0$ and $\alpha=\beta=\tfrac12$, and by \ref{P1}, the function
$\tau\mapsto\Phi(\abs\tau/2,\abs\tau/2)=\abs\tau^ph(0)$ is convex. It vanishes at $0$ and equals
$h(0)$ at $\pm1$, so $0\le h(0)$.

(b) For $t=0$ this is (a). Let $-1\le t<0$ and $u=\sqrt{-t}$. Then $0-u^2=t=x_t^2-y_t^2$, and
$u\le y_t$ because $y_t^2-u^2=(1+t)^2/4$. \Cref{lem:basic}\,(b), from $(0,u)$ to $(x_t,y_t)$,
and \ref{P1} give $(-t)^{p/2}h(-1)=u^p\Phi(0,1)=\Phi(0,u)\le h(t)$. The second inequality is the
first for the mirror image, whose function $h$ is $s\mapsto h(-s)$.

(c) Since $\Phi(0,0)=0$, the point $(x,y)$ is not the origin, and $h(\eta)<0$ for
$\eta=(x-y)/(x+y)\in[-1,1]$. If $\eta\le0$, (b) gives $(-\eta)^{p/2}h(-1)<0$, hence
$h(-1)<0$; if $\eta\ge0$, it gives $h(1)<0$.

(d) The supporting line at $s=-1$ and \cref{lem:ends} at $s=t$ give $h(t)-m(1+t)\le h(-1)$ and
$h(t)\ge h(-1)\bigl(1-\tfrac p2(1+t)\bigr)$. Hence $m(1+t)\ge-\tfrac p2h(-1)(1+t)$, which is the
first claim if $t>-1$. Likewise $h(t)+m(1-t)\le h(1)$ and
$h(t)\ge h(1)\bigl(1-\tfrac p2(1-t)\bigr)$ give $m(1-t)\le\tfrac p2h(1)(1-t)$.
\end{proof}

\begin{proof}[Proof of \cref{lem:hard}]
(a) If $p>2$, then $1-a=2/p$, and \cref{lem:ends} at $s=-a$ gives
$h(-a)\ge h(-1)(1-\tfrac p2(1-a))=0$. If $p<2$, then $1+a=2/p$, and
$h(-a)\ge h(1)(1-\tfrac p2(1+a))=0$.

(b) Let $p>2$, let $\lambda(s)=h(-1)(1-\tfrac p2(s+1))$ be the line of \cref{lem:ends} at $-1$,
with the slope $m_-=-\tfrac p2h(-1)$, and let $\ell(s)=H+m(s-t)$, $H=h(t)$. Then
$\lambda(-a)=0$, $H\ge\lambda(t)$ by \cref{lem:ends}, and $m\ge m_-$ by \cref{lem:orient}\,(d).
The identity
\[
  p\,\ell(-a)=p(H-mt)-(p-2)m=p\bigl(H-\lambda(t)\bigr)+(m-m_-)\bigl(-(pt+p-2)\bigr)
\]
shows that $\ell(-a)\ge0$ if $pt+p-2\le0$, that is, if $t\le-a$. So a hard base has $t>-a$. For
$p<2$ the same argument with the line $h(1)(1+\tfrac p2(s-1))$ at $1$, whose slope is at least
$m$, shows that $\ell(-a)\ge0$ if $t\ge-a$.

(c) By \cref{lem:orient}\,(d), $m>0$. If $p>2$, then $t>-a$ by (b), so
$1-\tfrac p2(t+1)<0$ and $h(t)\ge h(-1)(1-\tfrac p2(t+1))>0$ by \cref{lem:ends}. If $p<2$, then
$t<-a$, and $h(t)=\ell(-a)-m(-a-t)<0$.
\end{proof}

\begin{proof}[Proof of \cref{lem:easy}]
(a) In conformal coordinates, $G_{c_0,m}=\sigma^{p-1}(c_0\sigma+m\delta)$ and
$B_p=\sigma^{p-1}(\abs{1-2/p}\,\sigma+\delta)$, so that
$G_{c_0,m}-mB_p=(c_0-\abs{1-2/p}\,m)\,\sigma^p$.

(b) For $p>2$, $c_0-\abs{1-2/p}\,m=(pc_0-(p-2)m)/p\ge0$, so (a) writes $G_{c_0,m}$ as a
nonnegative combination of $B_p$, which is quasiconvex by \cref{cor:burkholder}, and of
$\abs A^p$, which is convex, hence polyconvex and quasiconvex (\cref{lem:supports}). Both are
continuous, and \cref{lem:qc-operations}\,(b) applies.

(c) For $1<p<2$, the identity of \cref{thm:low} writes $G_{c_0,m}$ as a combination of $B_p$
and of $A\mapsto B_p(JA)$ with nonnegative coefficients, and both are quasiconvex and continuous
(\cref{cor:burkholder}).
\end{proof}

\subsection{Restricted supports}

\begin{lemma}\label{lem:psi-low}
Let $1<p<2$, $a_0<0$, $a_2<0$ and $a_1\in\R$, and let $\psi(l)=a_0l^p+a_1l^{p-1}+a_2l^{p-2}$
for $l>0$. There is $l>0$ with $pa_0l^2+(p-1)a_1l+(p-2)a_2=0$ and $\psi(l)\ge\psi(1)$.
\end{lemma}

\begin{proof}
Let $K_0=\abs{\psi(1)}+1$. For $l_1\ge\max\{1,(a_1+K_0)/(-a_0)\}$,
$\psi(l_1)\le l_1^{p-1}(a_0l_1+a_1)\le-K_0$, because $a_2<0$. For $l_0\le1$ with
$l_0^{p-2}\ge(\abs{a_1}+K_0)/(-a_2)$, which exists because $p<2$,
$\psi(l_0)\le\abs{a_1}+a_2l_0^{p-2}\le-K_0$. So the maximum of the continuous function $\psi$
on $[l_0,l_1]$ is attained at an inner point $l$, where $\psi(l)\ge\psi(1)$ and
$0=\psi'(l)=l^{p-3}(pa_0l^2+(p-1)a_1l+(p-2)a_2)$.
\end{proof}

\begin{lemma}\label{lem:decreasing}
Let $p,a_0,a_1,a_2\in\R$, let $0<s\le1$, and let $pa_0l^2+(p-1)a_1l+(p-2)a_2\le0$ for
$l\in[s,1]$. Then $a_0+a_1+a_2\le a_0s^p+a_1s^{p-1}+a_2s^{p-2}$.
\end{lemma}

\begin{proof}
The function $\psi(l)=a_0l^p+a_1l^{p-1}+a_2l^{p-2}$ has the derivative
$l^{p-3}(pa_0l^2+(p-1)a_1l+(p-2)a_2)\le0$ on $[s,1]$, so $\psi(1)\le\psi(s)$.
\end{proof}

The next two lemmas are the mirror images of \cref{lem:caseI-algebra,lem:mu,lem:line}.

\begin{lemma}\label{lem:caseII-algebra}
Let $p,H,m,t,c,v\in\R$ with $t<1$, $c(1-t)=(1+t)m-pH$ and $v(1-t)=2C$, where $C=pH-2tm$, and
let $\alpha=ct-(p-1)H$. Then
\begin{align}
  &v\,x_t+2ct=pH,\label{eq:II1}\\
  &v+c(1+t)=pH+(1-t)m,\label{eq:II2}\\
  &p\alpha+(p-1)v\,x_t+(p-2)ct=0.\label{eq:II3}
\end{align}
If moreover $p<2$, $t>-1$, $H<0$ and $C\ge0$, then $\alpha<0$. If $1<p<2$, $c>0$, $v\ge0$ and
$p(H-mt)+(p-2)m<0$, then $\mu=(p-1)v/((2-p)c)$ satisfies $0\le\mu\le2$.
\end{lemma}

\begin{proof}
The three relations are computations from the definitions, and so are
\[
  2(1-t)\,\alpha=-\bigl((1+t)C+(p-2)(1-t)H\bigr),\qquad
  (2-\mu)(2-p)c(1-t)=-2\bigl(p(H-mt)+(p-2)m\bigr).
\]
Under the respective hypotheses the first right side is negative and the second positive.
Finally $\mu\ge0$ because $v\ge0$ and $c>0$.
\end{proof}

\begin{lemma}\label{lem:line-low}
Let $\Phi$ be a profile of degree $p<2$. Let $\abs t<1$, $H,m\in\R$, and suppose that
$G_{H-mt,m}\le\Phi$ on the quadrant and $\Phi(x_t,y_t)=H$. Let $c>0$ and $v$ satisfy
$c(1-t)=(1+t)m-pH$ and $v(1-t)=2C$, let $\alpha=ct-(p-1)H$, and let $\mu(2-p)c=(p-1)v$ with
$2t\le\mu\le2$. If $X,Y\ge0$ satisfy
\begin{equation}\label{eq:stationary-low}
  p\alpha+(p-1)v\,X+(p-2)c\,(X^2-Y^2)=0,
\end{equation}
then $\alpha+vX+c(X^2-Y^2)\le\Phi(X,Y)$.
\end{lemma}

\begin{proof}
Subtracting \eqref{eq:II3} from \eqref{eq:stationary-low} and dividing by $(2-p)c$ gives
$X^2-Y^2=t+\mu\tau$ with $\tau=X-x_t$. Since $2t\le\mu\le2$, there is $\theta$ with
$\cos\theta\,(1-t)=1+t-\mu$. Consider the rank-one line with $\alpha'=1$ and
$\beta'=e^{i\theta}$ in \cref{prop:tangent}. Then $\abs{x_t+\tau}=X$, and, since $2y_t=1-t$ and
$x_t^2-y_t^2=t$,
\[
  \abs{y_t+\tau e^{i\theta}}^2=y_t^2+\tau(1+t-\mu)+\tau^2=(x_t+\tau)^2-t-\mu\tau=Y^2 .
\]
By \eqref{eq:II2}, $\Phi_x=v+c(1+t)$, and by the definition of $c$, $\Phi_y=-c(1-t)$. So
\cref{prop:tangent} gives
\[
  \Phi(X,Y)\ge H+\tau(\Phi_x+\Phi_y\cos\theta)=H+\tau(v+c\mu)=\alpha+vX+c(X^2-Y^2),
\]
where the last equality uses $\alpha+vx_t+ct=H$, a consequence of \eqref{eq:II1}, and
$X^2-Y^2=t+\mu\tau$.
\end{proof}

\begin{proof}[Proof of \cref{lem:rs-identities}]
(b) $(1+et)(c-ev/2)=epH+(1-et)m-e(pH-2tm)=(1+et)m$, and $\alpha+v/2=H-tm$ is a computation.

(a) If $x+y=1$, then $r=(1-e\delta)/2$ and $x^2-y^2=\delta$, so by (b)
$q(x,y)=(\alpha+v/2)+\delta(c-ev/2)=H-tm+m\delta$. At $(x_t,y_t)$, $\delta=t$ and this is $H$.

(c) The point $(x/s,y/s)$ is on the segment, with $\delta/s$ in place of $\delta$. Since $r$ and
$x^2-y^2$ are homogeneous of degrees $1$ and $2$, (a) gives
$\alpha s^p+vrs^{p-1}+c(x^2-y^2)s^{p-2}=s^pq(x/s,y/s)=s^{p-1}((H-mt)s+m\delta)
=G_{H-mt,m}(x,y)$.

(d) Here $1\pm et>0$, and $v\ge0$ because $C\ge0$ (\cref{cor:BE}). Let $h(-1)<0$. By
\cref{lem:hard}\,(c), $m>0$ and $eH>0$, so the numerator $p\,eH+(1-et)m$ of $c$ is positive.
Finally, a computation gives
\[
  2(1+et)\,\alpha=-\bigl((1-et)C+(p-2)(1+et)H\bigr),
\]
and $(p-2)H=\abs{p-2}\,eH>0$; so $\alpha<0$.
\end{proof}

\begin{proof}[Proof of \cref{prop:restricted}]
Let $C=pH-2tm\ge0$. By \cref{lem:hard}\,(c) and \cref{lem:rs-identities}\,(d), $m>0$, $eH>0$,
$c>0$, $v\ge0$ and $\alpha<0$. For a point $(x,y)$ let
$\psi(l)=\alpha l^p+vrl^{p-1}+c(x^2-y^2)l^{p-2}$, where $r$ is taken at $(x,y)$, so that
$\psi(1)=q(x,y)$ and $\psi(l)=l^pq(x/l,y/l)$.

(a) Let $p<2$, so that $e=-1$, $r=x$, $c(1-t)=(1+t)m-pH$ and $v(1-t)=2C$. By
\cref{lem:hard}, $t<-a<0$. Let $x<y$; then $c(x^2-y^2)<0$. \Cref{lem:psi-low} with $a_0=\alpha$,
$a_1=vx$ and $a_2=c(x^2-y^2)$ gives $l>0$ with $\psi(l)\ge\psi(1)$ and
$p\alpha l^2+(p-1)vxl+(p-2)c(x^2-y^2)=0$. Dividing by $l^2$, the point $(X,Y)=(x/l,y/l)$
satisfies \eqref{eq:stationary-low}. By \cref{lem:caseII-algebra}, $\mu=(p-1)v/((2-p)c)$ lies in
$[0,2]$, and $2t<0$. So \cref{lem:line-low} gives $q(X,Y)\le\Phi(X,Y)$, and
\[
  \Phi(x,y)=l^p\Phi(X,Y)\ge l^pq(X,Y)=\psi(l)\ge\psi(1)=q(x,y).
\]

(b) Let $p>2$, so that $e=1$, $r=y$, and $c$, $v$, $\alpha$ are the numbers $c$, $\gamma$, $k$
of \cref{sec:caseI}. The hard base satisfies \eqref{eq:caseIhyp} and $H>0$. The proofs of
\cref{lem:caseI-algebra,lem:mu,lem:line} use $H>0$, \eqref{eq:caseIhyp} and
\cref{prop:tangent}, but not the sign of $\Phi$, so these lemmas hold here. Let
$\Phi(x,y)\ge0$. If $q(x,y)\le0$, then $q(x,y)\le\Phi(x,y)$. Otherwise $\psi(1)>0$, and the
argument after \cref{lem:line} gives $\psi(1)\le\Phi(x,y)$.

(c) At the origin $q=\alpha<0=\Phi(0,0)$. Let $0<s=x+y\le1$. By \cref{lem:decreasing} and
\cref{lem:rs-identities}\,(c), it is enough to show that
$P(l)=p\alpha l^2+(p-1)vrl+(p-2)c(x^2-y^2)\le0$ for $l\in[s,1]$: then
$q(x,y)=\psi(1)\le\psi(s)=G_{H-mt,m}(x,y)\le\Phi(x,y)$.

Let $p<2$. Where $x<y$ the claim is (a). Let $x\ge y$, so that $\delta\ge0\ge ts$. With $\mu$
as in (a), a computation from \eqref{eq:II3} gives
\[
  P(l)=(2-p)c\Bigl(s(ts-\delta)\Bigl(1-\frac\mu2\Bigr)+(l-s)\Bigl(t\Bigl(1-\frac\mu2\Bigr)(l+s)
  +\frac\mu2(\delta-l)\Bigr)\Bigr).
\]
For $l\ge s\ge\delta$ each term in the bracket is at most $0$, since $t\le0$ and $0\le\mu\le2$.

Let $p>2$. Where $\Phi\ge0$ the claim is (b). Let $\Phi(x,y)<0$. Then
$G_{H-mt,m}(x,y)<0$, that is, $Hs+m(\delta-ts)<0$, and since $H>0$ and $m>0$, $\delta<ts$.
Let $b=(p-1)v/(2(p-2)c)\ge0$. Computations from \eqref{eq:I1} and the definitions give
\begin{align*}
  (1-b)\cdot2(p-2)c(1+t)&=-2\bigl(p(H-mt)-(p-2)m\bigr)>0,\\
  c\bigl((p-2)b(1-t)+2(p-1)t\bigr)&=(p-1)pH>0 .
\end{align*}
So $b<1$, and $(p-1)(b(1-t)+2t)>b(1-t)\ge0$. From $b(1-t)+2t>0$ and the identity
$b(1-t)+2t=(b+t)(1-t)+t(1+t)$ it follows that $b+t\ge0$ and $\beta=b(1-t)+t\ge0$. Now
$P(l)=(p-2)c\,Q(l)$ with $Q(l)=-\beta l^2+2byl+x^2-y^2$, and
\[
  (1+t)s\,Q(l)=(ts-\delta)\bigl(-(1-b)(b+t)l^2-(bl-s)^2\bigr)+2x(l-s)(-\beta l-ts).
\]
For $l\ge s$ the first term is at most $0$, and so is the second, because
$-\beta l-ts\le-(\beta+t)s=-(b(1-t)+2t)s<0$. Hence $Q(l)\le0$.
\end{proof}

\begin{proof}[Proof of \cref{prop:monotone}]
Write $q_i$ and $\alpha_i$ for the function $q$ and the number $\alpha$ of the base $i$, and
$P_i=S_i(x_{t_i},y_{t_i})$; at $P_i$, $r=r_i$ and $x^2-y^2=d_i$. By homogeneity and
\cref{lem:rs-identities}\,(a),
\[
  S_i^p\alpha_i+\hat v_ir_i+\hat c_id_i=S_i^pq_i(x_{t_i},y_{t_i})=S_i^ph(t_i)=\Phi(P_i).
\]
Let $j\ne i$. Then $S_i^p\alpha_i+\hat v_ir_j+\hat c_id_j=S_i^pq_i(P_j/S_i)$. If $p<2$, the point
$P_j/S_i$ has $x<y$, because $t_j<-a<0$ by \cref{lem:hard}\,(b); if $p>2$,
$\Phi(P_j/S_i)=(S_j/S_i)^ph(t_j)>0$ by \cref{lem:hard}\,(c). By \cref{prop:restricted}\,(a)
or~(b), $S_i^pq_i(P_j/S_i)\le S_i^p\Phi(P_j/S_i)=\Phi(P_j)$. Adding the two inequalities, for
$(i,j)=(1,2)$ and $(2,1)$, and subtracting the two equalities gives the claim.
\end{proof}

\subsection{Radial integration}

\begin{proof}[Proof of \cref{lem:radial}]
(a) $\int_0^\sigma(n_0+n_1r)\,dr=n_0\sigma+n_1\sigma^2/2$, so
$\Psi^e_N(X)=(\sigma+e\delta)(n_0+n_1\sigma)-2n_0\sigma-n_1\sigma^2=-n_0(\sigma-e\delta)
+en_1\sigma\delta$, and $\sigma\delta=\det X$ by \cref{lem:coordinates}\,(a).

(b) $\int_0^\sigma(r-\rho)_+\,dr=(\sigma-\rho)_+^2/2$, so
$\Psi^e_N(X)=(\sigma-\rho)_+\bigl(\sigma+e\delta-(\sigma-\rho)\bigr)=(\sigma-\rho)_+(\rho+e\delta)$.
The matrix $\rho^{-1}X$ has the numbers $\sigma/\rho$ and $\delta/\rho$, so
$\rho^2K_e(\rho^{-1}X)=\rho^2(\sigma/\rho-1)_+(1+e\delta/\rho)$ is the same.

(c) By \cref{lem:coordinates}\,(d), $JX$ has the coordinates $(\bar w,\bar z)$. If $\sigma\le1$,
then $L(JX)+\det X=(\abs w^2-\abs z^2)+(\abs z^2-\abs w^2)=0=K_1(X)$, and likewise
$L(X)-\det X=0=K_{-1}(X)$. If $\sigma>1$, then, as $2\abs w=\sigma-\delta$ and
$2\abs z=\sigma+\delta$,
\[
  L(JX)+\det X=2\abs w-1+\sigma\delta=(\sigma-1)(1+\delta),\qquad
  L(X)-\det X=2\abs z-1-\sigma\delta=(\sigma-1)(1-\delta).
\]
The truncated determinant $L$ is continuous, and it is quasiconvex
by~\cite[Theorem~1]{SharpHI}, as formalized in~\cite{IwaniecLean} for the test maps of this
paper (\cref{sec:intro-low}). So is $X\mapsto L(JX)$ (\cref{lem:reflect}). The functions
$\pm\det X$ are polyaffine, hence quasiconvex (\cref{lem:supports}). By
\cref{lem:qc-operations}\,(b), $K_1$ and $K_{-1}$ are quasiconvex.

(d) $\Psi^e_N(X)-\Psi^e_{N'}(X)=(\sigma+e\delta)(N-N')(\sigma)-2\int_0^\sigma(N-N')$. Since
$\abs{\sigma+e\delta}\le2\sigma$ and $\sigma\le R$, each term is at most $2R\varepsilon$ in
absolute value.
\end{proof}

\begin{proof}[Proof of \cref{lem:hinge}]
Let $\eta>0$ be such that $\abs{N(r)-N(r')}<\varepsilon$ for $r,r'\in[0,R]$ with
$\abs{r-r'}\le\eta$, and let $n$ be an integer with $d=R/n<\eta$. Put $\rho_j=jd$ and
$s_j=(N(\rho_{j+1})-N(\rho_j))/d$ for $0\le j<n$. By convexity $s_0\le s_1\le\dots\le s_{n-1}$.
Let
\[
  P(r)=N(0)+s_0r+\sum_{j=1}^{n-1}(s_j-s_{j-1})(r-\rho_j)_+ .
\]
For $\rho_j\le r\le\rho_{j+1}$ the terms with an index above $j$ vanish, and the telescoping sum
gives $P(r)=N(\rho_j)+s_j(r-\rho_j)$, the chord of $N$ over $[\rho_j,\rho_{j+1}]$. There
\[
  \bigl(N(r)-P(r)\bigr)d=(\rho_{j+1}-r)\bigl(N(r)-N(\rho_j)\bigr)+(r-\rho_j)\bigl(N(r)
  -N(\rho_{j+1})\bigr),
\]
which is less than $\varepsilon d$ in absolute value. Take $n_1=s_0$, $w_j=s_j-s_{j-1}\ge0$
and the nodes $\rho_j>0$, $1\le j<n$.
\end{proof}

\begin{proof}[Proof of \cref{prop:radial}]
\emph{Step 1.} Let $N(r)=n_0+n_1r+\sum_iw_i(r-\rho_i)_+$ with $n_0\le0$, $w_i\ge0$ and
$\rho_i>0$. The function $\Psi^e_N$ is linear in $N$, so it is the sum of
$\Psi^e_{n_0+n_1r}$ and of the functions $w_i\Psi^e_{(r-\rho_i)_+}$. By
\cref{lem:radial}\,(a),
$\Psi^e_{n_0+n_1r}(X)=-n_0(1-e)\abs z-n_0(1+e)\abs w+en_1\det X$, which is quasiconvex by
\cref{lem:qc-operations}\,(c), since $-n_0(1\pm e)\ge0$. By \cref{lem:radial}\,(b) and (c),
$\Psi^e_{(r-\rho_i)_+}(X)=\rho_i^2K_e(\rho_i^{-1}X)$ is quasiconvex, by
\cref{lem:qc-operations}\,(a) and~(b). These functions are continuous, and
\cref{lem:qc-operations}\,(b) gives the quasiconvexity of $\Psi^e_N$.

\emph{Step 2.} Let $N$ be as in the theorem, $A\in\M$, and $\varphi$ a test map with the
Lipschitz constant $K$ such that $\Psi^e_N(A+\nabla\varphi)$ is integrable. Let
$R=\abs A+4K+1$; then $\abs{A+\nabla\varphi(x)}\le R$ for every $x$. Let $\eta>0$ and
$\varepsilon=\eta/(8R)$. \Cref{lem:hinge} on $[0,R]$ gives a function $N'$ as in Step~1, with
$N'(0)=N(0)\le0$ and $\abs{N-N'}\le\varepsilon$ on $[0,R]$. By \cref{lem:radial}\,(d),
$\abs{\Psi^e_N-\Psi^e_{N'}}\le4R\varepsilon=\eta/2$ at $A$ and at the matrices
$A+\nabla\varphi(x)$. By Step~1,
\[
  \Psi^e_N(A)\le\Psi^e_{N'}(A)+\frac\eta2\le\int_{(0,1)^2}\Psi^e_{N'}(A+\nabla\varphi)\,dx
  +\frac\eta2\le\int_{(0,1)^2}\Psi^e_N(A+\nabla\varphi)\,dx+\eta .
\]
As $\eta>0$ is arbitrary, $\Psi^e_N$ is quasiconvex.
\end{proof}

\begin{proof}[Proof of \cref{cor:burkholder}]
(a) Let $p>0$ and $\sigma\ge0$. Then $\int_0^\sigma N_p=\sigma^p/p$, so that, for every $e\in\R$,
\[
  \Psi^e_{N_p}(X)=(\sigma+e\delta)\sigma^{p-1}-\frac2p\sigma^p
  =\sigma^{p-1}\Bigl(\Bigl(1-\frac2p\Bigr)\sigma+e\delta\Bigr),
\]
and $\Psi^e_{-N_p}=-\Psi^e_{N_p}$. In conformal coordinates
$B_p(X)=\sigma^{p-1}(\abs{1-2/p}\,\sigma+\delta)$, and $B_p(JX)$ is the same with $-\delta$ in
place of $\delta$, by \cref{lem:coordinates}\,(d). For $p\ge2$, $\abs{1-2/p}=1-2/p$, which gives
the first two identities; for $p\le2$, $\abs{1-2/p}=2/p-1$, which gives the other two.

(b) The function $N_p$ is continuous for $p\ge1$, and $N_p(0)=0$ for $p>1$. For $p\ge2$ it is
convex on $[0,\infty)$, since $p-1\ge1$, and \cref{prop:radial} with $e=1$ and (a) show that
$B_p$ is quasiconvex. For $1<p\le2$ the function $r\mapsto r^{p-1}$ is concave on $[0,\infty)$,
so $-N_p$ is convex there, and \cref{prop:radial} with $e=-1$ and (a) show that $B_p$ is
quasiconvex. The mirror image is quasiconvex by \cref{lem:reflect}.
\end{proof}

\subsection{The path}

\begin{lemma}\label{lem:legendre}
Let $u_0\le u_1$, and let $h_0$ be strictly convex and continuous on $[u_0,u_1]$. There are
functions $T^\ast,k^\ast\colon\R\to\R$ such that, for every $m\in\R$:
\begin{enumerate}[label=\textup{(\roman*)}]
\item $T^\ast(m)\in[u_0,u_1]$, $k^\ast(m)=h_0(T^\ast(m))-mT^\ast(m)$, and
  $k^\ast(m)+ms\le h_0(s)$ for $s\in[u_0,u_1]$;
\item if $t\in[u_0,u_1]$ and $h_0(t)+m(s-t)\le h_0(s)$ for $s\in[u_0,u_1]$, then
  $T^\ast(m)=t$;
\item $k^\ast$ has the derivative $-T^\ast(m)$ at $m$.
\end{enumerate}
Moreover, $T^\ast$ is nondecreasing and continuous.
\end{lemma}

\begin{proof}
Let $T^\ast(m)$ be a point where the continuous function $t\mapsto h_0(t)-mt$ attains its
minimum on $[u_0,u_1]$, and $k^\ast(m)$ the minimum; this is (i). If $t$ is as in (ii), it is
also a minimizer, and if $t\ne T^\ast(m)$, strict convexity makes the midpoint smaller; so the
minimizer is unique, which is (ii). Adding the minimality inequalities for $m_1<m_2$ at
$T^\ast(m_2)$ and at $T^\ast(m_1)$ gives $(m_2-m_1)(T^\ast(m_2)-T^\ast(m_1))\ge0$, so
$T^\ast$ is nondecreasing. A one-sided limit of $T^\ast$ at $m$ is a minimizer for the slope
$m$, by the continuity of $h_0$, hence equals $T^\ast(m)$; so $T^\ast$ is continuous. Finally,
the minimality of $T^\ast(m')$ and of $T^\ast(m)$ give
\[
  -(m'-m)\bigl(T^\ast(m')-T^\ast(m)\bigr)\le k^\ast(m')-k^\ast(m)+(m'-m)T^\ast(m)\le0,
\]
and the left side is $o(\abs{m'-m})$ by the continuity of $T^\ast$. This is (iii).
\end{proof}

\begin{lemma}\label{lem:quadrature}
Let $b>0$ and $\nu>0$, and let $\rho$ be continuous on an open set $U\supseteq[0,b)$, with
$\rho(\xi)\ge\nu/(b-\xi)$ for $\xi\in[0,b)$. For every $R\ge1$ there is $\bar\xi\in[0,b)$ with
$\exp\int_0^{\bar\xi}\rho=R$. If $U$ contains the closed interval between $0$ and $\xi$, the
function $u\mapsto\exp\int_0^u\rho$ has the derivative $\rho(\xi)\exp\int_0^\xi\rho$ at $\xi$.
\end{lemma}

\begin{proof}
For $\xi\in[0,b)$, $\int_0^\xi\rho\ge\nu\log(b/(b-\xi))$, which tends to $\infty$ as
$\xi\to b$. The function $\xi\mapsto\exp\int_0^\xi\rho$ is continuous on $[0,b)$ and equals $1$
at $0$, so it takes the value $R$. The derivative is the fundamental theorem of calculus.
\end{proof}

\begin{proof}[Proof of \cref{prop:path}]
The function $h$ is continuous on $[-1,1]$, as the restriction of the convex function $\tilde h$
on $\R$ (\cref{lem:basic}\,(a)). Let $T^\ast$, $k^\ast$ be the functions of
\cref{lem:legendre} for $h$ on $[-1,1]$. The supporting line of slope $m_0$ at $t_0$ gives
$T^\ast(m_0)=t_0$. Let $m_a$ be the slope of a supporting line of $h$ at $-a$
(\cref{lem:minorant}), so that $T^\ast(m_a)=-a$, and let $\xi_a=e(m_0-m_a)$, so that
$m_0-e\xi_a=m_a$. By \cref{lem:hard}\,(b), $e(t_0+a)>0$; since $T^\ast$ is nondecreasing, this
forces $e(m_0-m_a)>0$, that is, $\xi_a>0$.

For $\xi\in\R$ let $m(\xi)=m_0-e\xi$, $T(\xi)=T^\ast(m(\xi))$, $k(\xi)=k^\ast(m(\xi))$ and
$g(\xi)=k(\xi)-a\,m(\xi)$. By \cref{lem:legendre}, $T$ is continuous, $k$ has the derivative
$eT(\xi)$, and $g$ has the derivative $e(T(\xi)+a)$. Since $pa=e(p-2)$, $p\,g(0)<0$ is the
condition of \cref{def:hard}, and $g(\xi_a)=h(-a)\ge0$ by \cref{lem:hard}\,(a). Let $b$ be the
smallest $\xi\in[0,\xi_a]$ with $g(\xi)\ge0$; then $b>0$ and $g<0$ on $[0,b)$. For
$\xi\in[0,\xi_a]$ the slope $m(\xi)$ lies between $m_0$ and $m_a$, so $T(\xi)$ lies between
$t_0$ and $-a$, and $\abs{T(\xi)}\le\vartheta=\max\{\abs{t_0},a\}<1$.

The function $\xi\mapsto2\xi-g(\xi)$ has the derivative $2-e(T+a)>0$, so that
$2\xi-g(\xi)\le2b-g(b)\le2b$ for $\xi\in[0,b)$, that is, $-g(\xi)\le2(b-\xi)$. Let
$\rho(\xi)=-(1+eT(\xi))/(p\,g(\xi))$, which is continuous on the open set where $g<0$. For
$\xi\in[0,b)$, $1+eT(\xi)\ge1-\vartheta$, so that $\rho(\xi)\ge(1-\vartheta)/(2p(b-\xi))$.
\Cref{lem:quadrature} gives $\bar\xi\in[0,b)$ with $S(\bar\xi)=R$ for
$S(\xi)=\exp\int_0^\xi\rho$; then $S(0)=1$, $S>0$, and $S$ has the derivative
$S\rho=-S(1+eT)/(pg)$ on $[0,\bar\xi]$. Finally, for $\xi\in[0,\bar\xi]$, $\abs{T(\xi)}<1$, the
line of slope $m(\xi)$ through $(T(\xi),h(T(\xi)))$ supports $h$ on $[-1,1]$ with the intercept
$k(\xi)$ (\cref{lem:legendre}\,(i)), and $g(\xi)<0$: $(T(\xi),m(\xi))$ is a hard base.
\end{proof}

\begin{lemma}\label{lem:self-improving}
Let $f_0$ be continuous on $[u_0,u_1]$, and let $\eta>0$. Suppose that for all
$\xi_1<\xi_2$ in $[u_0,u_1]$ with $\xi_2-\xi_1\le\eta$ and every $D$ with
$\abs{f_0(\zeta)-f_0(\xi_1)}\le D$ for $\zeta\in[\xi_1,\xi_2]$, we have
$f_0(\xi_1)-f_0(\xi_2)\le D/2$. Then $f_0$ is nondecreasing on $[u_0,u_1]$.
\end{lemma}

\begin{proof}
It is enough to compare points $\zeta_1<\zeta_2$ with $\zeta_2-\zeta_1\le\eta$, and to chain.
Suppose $f_0(\zeta_2)<f_0(\zeta_1)$. Let $\xi_2$ be a minimizer of $f_0$ on
$[\zeta_1,\zeta_2]$, and $\xi_1$ a maximizer of $f_0$ on $[\zeta_1,\xi_2]$. Then
$f_0(\xi_1)\ge f_0(\zeta_1)>f_0(\zeta_2)\ge f_0(\xi_2)$, so $\xi_1<\xi_2$, and
$D=f_0(\xi_1)-f_0(\xi_2)>0$ bounds $\abs{f_0-f_0(\xi_1)}$ on $[\xi_1,\xi_2]$. The hypothesis
gives $D\le D/2$, a contradiction.
\end{proof}

\begin{proof}[Proof of \cref{lem:path}]
On $[0,\bar\xi]$ we have $g<0$ and $1+eT>0$, so $S'>0$. We write $S$, $T$, $m$, $c$, $M$, $f$
for their values at $\xi$, and use $e^2=1$ and $H=h(T)=k+mT$.

(a) Since $m'=-e$, $M'=(p-1)S^{p-2}S'm-eS^{p-1}$, while $efS'=S^{p-2}cS'$. These agree if and
only if $S'(c-(p-1)m)=-eS$, that is, by the definition of $S'$, if and only if
\begin{equation}\label{eq:path-M}
  (1+eT)\bigl(c-(p-1)m\bigr)=e\,p\,g,
\end{equation}
which is a computation from the definition of $c$, with $pa=e(p-2)$.

(b) The function $f$ is continuous on $[0,\bar\xi]$, because $S$, $T$ and $k$ are. Let
$\xi_1<\xi_2$ in $[0,\bar\xi]$, and write $S_i$, $T_i$, $M_i$, $f_i$ for the values at $\xi_i$
and $r_i=S_i(1-eT_i)/2$. Apply \cref{prop:monotone} to the hard bases $(T_i,m(\xi_i))$ with the
scales $S_i$. By \cref{lem:rs-identities}\,(b), $v=2e(c-m)$, and a computation turns its
inequality into
\[
  (f_2-f_1)(S_2-S_1)\,W\ge2(r_2-r_1)\,Y,\qquad W=S_2+eS_1T_1,\quad
  Y=e(M_2-M_1)-f_1(S_2-S_1).
\]
Let $W_0=1-\max_{[0,\bar\xi]}\abs T>0$. Since $S$ is increasing with $S(0)=1$, $S_2>S_1\ge1$,
and $W\ge S_1(1-\abs{T_1})\ge W_0$. The function $S(1-eT)/2$ is uniformly continuous on
$[0,\bar\xi]$; let $\eta>0$ be such that $\abs{r_2-r_1}\le W_0/4$ when $\xi_2-\xi_1\le\eta$.
Let $\abs{f-f_1}\le D$ on $[\xi_1,\xi_2]$. By (a), the functions $DS\mp(eM-f_1S)$ have the
derivatives $(D\mp(f-f_1))S'\ge0$ on $[\xi_1,\xi_2]$, so that $\abs Y\le D(S_2-S_1)$. Hence
$(f_2-f_1)W\ge-DW_0/2$. If $f_1-f_2>D/2$, then $(f_2-f_1)W\le(f_2-f_1)W_0<-DW_0/2$, a
contradiction. So $f_1-f_2\le D/2$, and \cref{lem:self-improving} applies.

(c) By (a), the function $eM-f(\xi')S$ has the derivative $(f-f(\xi'))S'$, which is at least $0$
on $[\xi',\bar\xi]$ and at most $0$ on $[0,\xi']$. So it is smallest at $\xi'$:
$eM(\xi)-f(\xi')S(\xi)\ge eM(\xi')-f(\xi')S(\xi')$.
\end{proof}

\begin{proof}[Proof of \cref{lem:envelope}]
(a) The lines depend continuously on $\xi\in[0,\bar\xi]$, so the supremum is finite at every
$s$, and a supremum of affine functions is convex.

(b) The line with the index $\xi$ gives $N(S(\xi))\ge eM(\xi)$, and \cref{lem:path}\,(c) gives
the converse inequality.

(c) Let $s\le1$ and $\xi\in[0,\bar\xi]$. By \cref{lem:path}\,(c) at the point $S(0)=1$,
$eM(\xi)+f(\xi)(1-S(\xi))\le eM(0)=em_0$, and $f(\xi)\ge f(0)=ec(0)$. Hence
\[
  eM(\xi)+f(\xi)(s-S(\xi))\le em_0+f(\xi)(s-1)\le em_0+ec(0)(s-1),
\]
which is the line with the index $0$. So $N(s)=em_0+ec(0)(s-1)$, and this is
$ec(0)s-v(0)/2$ by \cref{lem:rs-identities}\,(b).

(d) $N$ is continuous, being convex on $\R$. Let
$J(\xi)=\alpha(0)+S\,eM-2\int_0^SN(r)\,dr-S^pk$. By (b), the fundamental theorem of calculus,
\cref{def:path}\,(ii) and \cref{lem:path}\,(a),
\[
  J'=e(SM'-MS')-pS^{p-1}S'k-eS^pT=S^{p-1}\bigl(S'(e(c-m)-pk)-eST\bigr).
\]
By the definition of $S'$, this vanishes because of the identity
\begin{equation}\label{eq:path-I}
  (1+eT)\bigl(pk-e(c-m)\bigr)=e\,p\,g\,T,
\end{equation}
a computation like \eqref{eq:path-M}. So $J$ is constant on $[0,\bar\xi]$. At $0$, $S=1$,
$M=m_0$, $k(0)=h(t_0)-m_0t_0$ and $\int_0^1N=ec(0)/2-v(0)/2$ by (c), so that
$J(0)=\alpha(0)+v(0)/2-(h(t_0)-m_0t_0)+(em_0-ec(0)+v(0)/2)=0$ by \cref{lem:rs-identities}\,(b).
With (b), $I(S(\xi))=J(\xi)+S(\xi)^pk(\xi)=S(\xi)^pk(\xi)$.
\end{proof}

\begin{proof}[Proof of \cref{prop:minorant}]
Let $N$ be the upper envelope of \cref{sec:path}. By \cref{lem:path}\,(b), $f$ is nondecreasing,
so \cref{lem:envelope} applies: $N$ is convex, and $N(0)=-v(0)/2\le0$ by
\cref{lem:rs-identities}\,(d). Since $T(0)=t_0$ and $m(0)=m_0$, $\alpha(0)=\alpha_0$. Let
$x,y\ge0$ with $s=x+y\le R$.

If $s\le1$, then $N(r)=n_0+n_1r$ for $r\le1$, with $n_0=-v(0)/2$ and $n_1=ec(0)$, by
\cref{lem:envelope}\,(c). The computation of \cref{lem:radial}\,(a) gives
$\Psi^e_N(x,y)=-n_0(\sigma-e\delta)+en_1\sigma\delta=v(0)r+c(0)(x^2-y^2)$, so that
$\alpha_0+\Psi^e_N$ is the restricted support $q$ at $(t_0,m_0)$, which lies below $\Phi$ by
\cref{prop:restricted}\,(c).

If $1<s\le R$, there is $\xi\in[0,\bar\xi]$ with $S(\xi)=s$, since $S$ is continuous, $S(0)=1$
and $S(\bar\xi)=R$. By \cref{lem:envelope}\,(b) and (d),
\[
  \alpha_0+\Psi^e_N(x,y)=I(s)+eN(s)\,\delta=s^pk(\xi)+s^{p-1}m(\xi)\,\delta
  =G_{k(\xi),m(\xi)}(x,y)\le\Phi(x,y),
\]
because $(T(\xi),m(\xi))$ is a hard base with the intercept $k(\xi)$.

At $(x_{t_0},y_{t_0})$, $s=1$ and $\alpha_0+\Psi^e_N=q(x_{t_0},y_{t_0})=h(t_0)$ by
\cref{lem:rs-identities}\,(a).
\end{proof}

\begin{proof}[Proof of \cref{cor:minorant}]
\Cref{prop:path} gives a path of hard bases from $(t_0,m_0)$ to the scale $R$, and
\cref{prop:minorant} the function $N$.
\end{proof}

\subsection{The proof of \texorpdfstring{\cref{thm:A}}{Theorem A}}

\begin{proof}[Proof of \cref{lem:reflect}]
By \cref{lem:coordinates}\,(d), $JA$ has the coordinates $(\bar w,\bar z)$. So, if $F$ is
isotropic, $F(JA)=\Phi(\abs w,\abs z)$, which is isotropic by \cref{lem:coordinates}\,(c), and
the function $h$ of its profile has at $-1$ the value $\Phi(1,0)=h(1)$. Homogeneity passes
because $J(\lambda A)=\lambda JA$. If $B$ has rank one, so has $JB$, and
$F(J(A+\tau B))=F(JA+\tau JB)$; so rank-one convexity passes.
Finally, if $\varphi=(\varphi_1,\varphi_2)$ is a test map, so is $(\varphi_1,-\varphi_2)$, with
the gradient $J\nabla\varphi$; Jensen's inequality for $f$ at $JA$ with this test map is Jensen's
inequality for $A\mapsto f(JA)$ at $A$ with $\varphi$.
\end{proof}

\begin{proof}[Proof of \cref{lem:perturb}]
(a) $\Theta$ depends only on $\abs z$ and $\abs w$, so it is isotropic by
\cref{lem:coordinates}\,(c); the matrix $\lambda A$ has the coordinates $(\lambda z,\lambda w)$;
and $\diag(x+y,x-y)$ has the coordinates $(x,y)$. Along a line,
$\tau\mapsto\abs{z_0+\tau\alpha}^p$ is convex for $p\ge1$, as the composition of a convex
function with values in $[0,\infty)$ and the convex nondecreasing function $r\mapsto r^p$ on
$[0,\infty)$; likewise for $w$.

(b) For $p>1$, $r\mapsto r^p$ is strictly convex on $[0,\infty)$, so $t\mapsto((1+t)/2)^p$ is
strictly convex on $[-1,1]$, and $t\mapsto((1-t)/2)^p$ and $h_0$ are convex.
\end{proof}

\begin{proof}[Proof of \cref{prop:interior}]
Let $\Phi$ be the profile of $F$, $x_0=\abs{z_0}$, $y_0=\abs{w_0}$, $\sigma_0=x_0+y_0=\abs{A_0}>0$
and $t_0=(x_0-y_0)/\sigma_0$, so that $\abs{t_0}<1$ and $(x_0,y_0)=\sigma_0(x_{t_0},y_{t_0})$.
By \cref{lem:minorant} there is a supporting slope $m_0$ of $h$ at $t_0$, and $m_0>0$ by
\cref{lem:orient}\,(d). Let $H=h(t_0)$.

If $p(H-m_0t_0)-e(p-2)m_0\ge0$, the Burkholder-type function $G=G_{H-m_0t_0,m_0}$ satisfies
$G(A)\le\Phi(\abs z,\abs w)=F(A)$ for every $A$, and, as $\det A_0=\sigma_0^2t_0$,
$G(A_0)=\sigma_0^p(H-m_0t_0+m_0t_0)=\sigma_0^ph(t_0)=F(A_0)$. It is continuous, and
quasiconvex by \cref{lem:easy}\,(b) or (c): since $m_0>0$, the assumption reads
$(p-2)m_0\le p(H-m_0t_0)$ if $p>2$, and $(2-p)\abs{m_0}\le p(H-m_0t_0)$ if $p<2$. So $F$ is
quasiconvex at $A_0$ by the remark after \cref{lem:supports}.

Otherwise $(t_0,m_0)$ is a hard base. Let $\varphi$ be a test map with the Lipschitz constant
$K$, for which $F(A_0+\nabla\varphi)$ is integrable. Then
$\abs{A_0+\nabla\varphi(x)}\le\sigma_0+4K=R\sigma_0$ with $R=1+4K/\sigma_0\ge1$. Let $N$ be given
by \cref{cor:minorant} for this $R$, and let
\[
  G(X)=\sigma_0^p\bigl(\alpha_0+\Psi^e_N(\sigma_0^{-1}X)\bigr).
\]
It is quasiconvex by \cref{prop:radial} and \cref{lem:qc-operations}, and continuous. If
$\abs X\le R\sigma_0$, the point $(\abs z/\sigma_0,\abs w/\sigma_0)$ has coordinate sum at
most $R$, and by \ref{P1}
\[
  G(X)=\sigma_0^p\bigl(\alpha_0+\Psi^e_N(\abs z/\sigma_0,\abs w/\sigma_0)\bigr)\le
  \sigma_0^p\,\Phi(\abs z/\sigma_0,\abs w/\sigma_0)=F(X),
\]
with equality at $A_0$, where the point is $(x_{t_0},y_{t_0})$. Hence
\[
  F(A_0)=G(A_0)\le\int_{(0,1)^2}G(A_0+\nabla\varphi)\,dx\le\int_{(0,1)^2}F(A_0+\nabla\varphi)\,dx.
  \qedhere
\]
\end{proof}

\begin{proof}[Proof of \cref{lem:dense}]
Let $A_0\in\M$ with the coordinates $(z_0,w_0)$, and let $\varphi$ be a test map. Let $u=z_0$ if
$z_0\ne0$ and $u=1$ otherwise, $v=w_0$ if $w_0\ne0$ and $v=1$ otherwise, and let $B$ have the
coordinates $(u,v)$. For $\lambda_n=1/(n+1)$ the matrix $A_n=A_0+\lambda_nB$ has the
coordinates $z_0+\lambda_nu$ and $w_0+\lambda_nv$, which are nonzero: each is
$(1+\lambda_n)z_0$ or $\lambda_n$, and likewise for $w$. So
$F(A_n)\le\int_{(0,1)^2}F(A_n+\nabla\varphi)\,dx$. As $n\to\infty$, $F(A_n)\to F(A_0)$ by
continuity, and the integrals converge to $\int_{(0,1)^2}F(A_0+\nabla\varphi)\,dx$ by dominated
convergence: the matrices $A_n+\nabla\varphi(x)$ lie in a bounded set, on which $F$ is
bounded.
\end{proof}

\begin{proof}[Proof of \cref{prop:oriented}]
A rank-one convex function is continuous (proof of \cref{prop:small}), and so is $\Theta$. Let
$0<\varepsilon<-h(-1)$ and $F_\varepsilon=F+\varepsilon\Theta$. By \cref{lem:perturb},
$F_\varepsilon$ satisfies~\Hyp, its profile is $\Phi+\varepsilon(x^p+y^p)$, and the function $h$
of this profile, $h(t)+\varepsilon(((1+t)/2)^p+((1-t)/2)^p)$, is strictly convex on $[-1,1]$
and equals $h(-1)+\varepsilon<0$ at $-1$. By \cref{prop:interior,lem:dense}, $F_\varepsilon$ is
quasiconvex. Let $A\in\M$ and let $\varphi$ be a test map for which $F(A+\nabla\varphi)$ is
integrable; so is $\Theta(A+\nabla\varphi)$. Then
\[
  F(A)+\varepsilon\Theta(A)\le\int_{(0,1)^2}F(A+\nabla\varphi)\,dx
  +\varepsilon\int_{(0,1)^2}\Theta(A+\nabla\varphi)\,dx,
\]
and $\varepsilon\to0$ gives Jensen's inequality for $F$.
\end{proof}

\begin{proof}[Proof of \cref{thm:A}]
Let $F\ge0$. For $0<p<1$, $F=0$ by \cref{prop:small}. For $p=1$, $F$ is convex by
\cref{cor:low}\,(b), hence polyconvex and quasiconvex (\cref{lem:supports}). For $1<p<2$ this is
\cref{cor:low}\,(a). For $p\ge2$, $F$ is polyconvex by \cref{thm:main}, hence quasiconvex.

Let $F(A)<0$ for some $A$. Then $p>1$, since $F=0$ for $p<1$ and $F\ge0$ for $p=1$. If $p=2$,
$F$ is polyconvex by \cref{thm:iff}\,(b), hence quasiconvex. Let $p\ne2$. The profile $\Phi$ of
$F$ is negative at $(\abs z,\abs w)$ for the coordinates of $A$, so $h(-1)<0$ or $h(1)<0$ by
\cref{lem:orient}\,(c). In the first case \cref{prop:oriented} applies to $F$. In the second, it
applies to $A\mapsto F(JA)$ by \cref{lem:reflect}, and then $F(A)=F(J(JA))$ is quasiconvex by
the same lemma, as $J^2=I$.
\end{proof}

\section{The formal verification}\label{app:formal}

The formalization uses Lean~4~\cite{Lean4} and Mathlib~\cite{Mathlib20} and is in the directory
\lean{Morrey/HomogeneousIsotropic} of the repository~\cite{Companion}. We describe its
statements, their verification and dependencies, and the correspondence between the paper and
the Lean declarations.

\subsection{The statements}\label{app:formal-statements}

{\sloppy
The file \lean{palomar/HomogeneousIsotropic/Challenge.lean} imports only Mathlib and states the
results of \cref{sec:intro} and \cref{prop:radial} as theorems of the namespace
\lean{HomogeneousIsotropicChallenge}. They are \lean{theoremA}, \lean{theorem1_1},
\lean{theorem1_2}, \lean{theorem1_3_a}, \lean{theorem1_3_b}, \lean{corollary1_4},
\lean{theorem1_5}, \lean{theorem1_5_identity}, \lean{corollary1_6_a}, \lean{corollary1_6_b},
\lean{proposition1_7} and \lean{proposition6_8}, together with the two identities of
\cref{lem:coordinates}\,(a), \lean{det_eq_norm_sq_sub} and \lean{opNorm_eq}. Its definitions are
the following.\par}
{\sloppy
\begin{itemize}
\item Matrices are \leanx{Matrix (Fin 2) (Fin 2) $\R$}, and the gradient of
  $\varphi\colon\R^2\to\R^2$ is the matrix \leanx{jac $\varphi$ x} of the derivatives of its
  components.
\item \leanx{IsRankOneConvex F}: $\tau\mapsto F(A+\tau B)$ is convex on $\R$ for every $A$ and
  every $B$ with \leanx{B.rank~=~1}.
\item \leanx{IsQuasiconvex F}: $F(A)\le\int_{(0,1)^2}F(A+\nabla\varphi)$ for every $A$ and every
  Lipschitz $\varphi$ vanishing outside the open unit square, for which $F(A+\nabla\varphi)$ is
  integrable.
\item \leanx{IsIsotropic F}: invariance under
  \leanx{Matrix.specialOrthogonalGroup (Fin 2) $\R$} on both sides. \leanx{IsPosHomogeneous p F}:
  $F(cA)=c^pF(A)$ for $c>0$, with the real power.
\item \leanx{Polyconvex F}: $F(A)=g(A,\det A)$ for a convex $g$ on the product of the matrices
  and $\R$. \leanx{HasPolyaffineSupportAt F A$_0$}: a function
  $k+\sum_{i,j}K_{ij}A_{ij}+c\det A$ lies below $F$ and equals $F$ at $A_0$.
\item \leanx{opNorm A} is the operator norm $\abs A$ of $A$ on Euclidean $\R^2$, and
  \lean{conformalPart}, \lean{anticonformalPart} are the coordinates $z$, $w$.
\item \leanx{burkholderType p c$_0$ m A} is $G_{c_0,m}(A)=\abs A^{p-2}(c_0\abs A^2+m\det A)$,
  and \leanx{burkholder p A} is $B_p(A)=(\abs{1-2/p}\abs A^2+\det A)\abs A^{p-2}$. With the real
  power of Lean, both are $0$ at $A=0$.
\item \leanx{multiplier p H m t} is the number $c^\ast$ of \cref{thm:closed}.
\item \leanx{radial e N X} is the function $\Psi^e_N(X)$ of \cref{prop:radial}, written with the
  moduli of \lean{conformalPart} and \lean{anticonformalPart}, and
  \leanx{radialProfile e N x y} is the same function of $x=\abs z$ and $y=\abs w$.
\end{itemize}\par}
The definitions are copied verbatim from the library, so the library's theorems have the
statements given in the Challenge file. The facts quoted from the literature in
\cref{sec:intro,sec:discussion}, such as the rank-one convexity of $B_p$ and the thresholds for
the integrands $f_\gamma$, are not formalized.

\subsection{Verification and dependencies}\label{app:formal-check}

{\sloppy
The theorems are proved in \lean{palomar/HomogeneousIsotropic/Solution.lean} from the library.
Lean's kernel checks the proofs. Every theorem depends only on the standard axioms \lean{propext},
\lean{Classical.choice} and \lean{Quot.sound}, as verified by an audit of every declaration of
the library. Lean's Comparator, with the configuration
\lean{palomar/HomogeneousIsotropic/comparator.json}, checks that the Solution proves the
statements of the Challenge with the same definitions, under Lean's kernel and the independent
kernel nanoda. Lean and Mathlib are pinned at version \texttt{v4.36.0-rc2}.\par}

\begin{sloppypar}
The library's proofs use the rank-one convexity and the quasiconvexity of the library
\lean{CalculusOfVariations}~\cite{CoVLean}, which are the notions above
(\lean{Morrey/HomogeneousIsotropic/Bridges.lean}). From that library they take two facts: minors
are null Lagrangians, which gives \cref{lem:supports}\,(a), and rank-one convex functions are
locally Lipschitz, which gives the continuity used in \cref{prop:small,prop:oriented}. From the
formalization~\cite{IwaniecLean} of~\cite{SharpHI} they take three theorems:
\end{sloppypar}
\begin{itemize}
\item the quasiconvexity of the truncated determinant (\lean{Agazzi.quasiconvex_integrandL}),
  for \cref{lem:radial}\,(c), hence for \cref{prop:radial}, and through it for
  \cref{cor:burkholder}, \cref{cor:low}\,(a) and \cref{thm:A};
\item the identity $\abs A=\abs z+\abs w$ for the operator norm (\lean{Agazzi.opNorm_eq});
\item the invariance of quasiconvexity under positive multiples
  (\lean{Agazzi.quasiconvex_const_mul}).
\end{itemize}
\begin{sloppypar}
\noindent With the first comes the continuity of the truncated determinant
(\lean{Agazzi.continuous_integrandL}). That formalization proves the quasiconvexity of $L$ for the
Lipschitz test maps vanishing outside the unit square (\cref{sec:intro-low}). The proofs of
\cref{sec:signed} are in the directory \lean{Morrey/HomogeneousIsotropic/Signed}: the truncated
determinant enters in its file \lean{Radial.lean}, and \cref{cor:burkholder} is in
\lean{BurkholderRadial.lean}. The identity $\abs A=\abs z+\abs w$ enters in
\lean{Morrey/HomogeneousIsotropic/Burkholder.lean}.
\Cref{lem:convex,lem:sublinear} are in the file
\lean{Morrey/Common/Basic/Analysis/Convex/Tangent.lean}, and the general lemmas of
\cref{app:signed} in the files \lean{Convex/Hinge.lean} and \lean{Convex/Legendre.lean} of the
same directory, \lean{Morrey/Common/Basic/MeasureTheory/Integral/IntervalIntegral/Primitive.lean}
and \lean{Morrey/Common/Basic/Topology/Order/SelfImproving.lean}. The hypotheses of
\cref{thm:main} are checked on an integrand that is not convex, $F(A)=(\det A)^2$ with $p=4$
(\lean{polyconvex_detSq}, \lean{not_convexOn_detSq}).
\end{sloppypar}

\subsection{Dictionary}\label{app:formal-dictionary}

{\sloppy
\Cref{tab:formal-dictionary} lists the Lean declarations of the statements of the paper. Names
are relative to the namespace \lean{HomogeneousIsotropic}; the prefix
\lean{HomogeneousIsotropicChallenge} marks the statements of the Challenge file. The general
lemmas, with the prefixes \lean{ConvexOn}, \lean{StrictConvexOn} and \lean{intervalIntegral},
and \lean{monotoneOn_of_forall_sub_le_half}, are not in this namespace and are named in full.
When a row has several declarations, they state the parts of the result: for a support, the
inequality and the equality at the base point.\par}

{\small
\begin{longtable}{@{}>{\raggedright\arraybackslash}p{0.25\textwidth}%
>{\raggedright\arraybackslash}p{0.71\textwidth}@{}}
\caption{The Lean declarations of the statements of the paper.}\label{tab:formal-dictionary}\\
\toprule
Result & Lean declarations\\
\midrule
\endfirsthead
\toprule
Result & Lean declarations\\
\midrule
\endhead
\bottomrule
\endfoot
\multicolumn{2}{@{}l}{\emph{\Cref{sec:intro}}}\\*[2pt]
\Cref{thm:A} & \lean{HomogeneousIsotropicChallenge.theoremA};
  \lean{quasiconvex_of_rankOneConvex_of_isotropic}, \lean{quasiconvex_of_rankOneConvex_of_pos}
  \\\addlinespace[2pt]
\Cref{thm:main} & \lean{HomogeneousIsotropicChallenge.theorem1_1};
  \lean{hasPolyaffineSupportAt_of_rankOneConvex}, \lean{polyconvex_of_rankOneConvex},
  \lean{quasiconvex_of_rankOneConvex}\\\addlinespace[2pt]
\Cref{thm:closed} & \lean{HomogeneousIsotropicChallenge.theorem1_2};
  \lean{polyaffine_multiplier_le}\\\addlinespace[2pt]
\Cref{thm:iff} & \lean{HomogeneousIsotropicChallenge.theorem1_3_a},
  \lean{HomogeneousIsotropicChallenge.theorem1_3_b}; \lean{polyconvex_iff_of_two_lt},
  \lean{nonneg_of_polyconvex}, \lean{polyconvex_iff_two}\\\addlinespace[2pt]
\Cref{cor:equiv} & \lean{HomogeneousIsotropicChallenge.corollary1_4};
  \lean{tfae_of_two_le}\\\addlinespace[2pt]
\Cref{thm:low} & \lean{HomogeneousIsotropicChallenge.theorem1_5},
  \lean{HomogeneousIsotropicChallenge.theorem1_5_identity}; \lean{exists_burkholderType_le},
  \lean{burkholderType_eq}\\\addlinespace[2pt]
$B_p$ in coordinates & \lean{burkholder_eq_coord}, \lean{burkholder_eq_burkholderTypeCoord}
  \\\addlinespace[2pt]
\Cref{cor:low}\,(a) & \lean{HomogeneousIsotropicChallenge.corollary1_6_a};
  \lean{quasiconvex_burkholder}, \lean{quasiconvex_of_rankOneConvex_of_burkholder},
  \lean{quasiconvex_of_rankOneConvex_of_lt_two}\\\addlinespace[2pt]
\Cref{cor:low}\,(b) & \lean{HomogeneousIsotropicChallenge.corollary1_6_b};
  \lean{nonneg_of_rankOneConvex_one}, \lean{convexOn_of_rankOneConvex_one}\\\addlinespace[2pt]
\Cref{prop:small} & \lean{HomogeneousIsotropicChallenge.proposition1_7};
  \lean{eq_zero_of_rankOneConvex_of_lt_one}\\[4pt]
\multicolumn{2}{@{}l}{\emph{\Cref{sec:setting,app:setting}}}\\*[2pt]
\Cref{lem:coordinates}\,(a) & \lean{HomogeneousIsotropicChallenge.det_eq_norm_sq_sub},
  \lean{HomogeneousIsotropicChallenge.opNorm_eq}; \lean{det_eq_norm_sq_sub},
  \lean{opNorm_eq}\\\addlinespace[2pt]
\Cref{lem:coordinates}\,(b) & \lean{rank_eq_one_iff_norm_eq}\\\addlinespace[2pt]
\Cref{lem:coordinates}\,(c) & \lean{ofComplex_mem_specialOrthogonalGroup},
  \lean{exists_ofComplex_eq_of_mem_specialOrthogonalGroup}, \lean{ofComplex_mul_mul}
  \\\addlinespace[2pt]
\Cref{lem:coordinates}\,(d) & \lean{conformalPart_reflection_mul},
  \lean{anticonformalPart_reflection_mul}\\\addlinespace[2pt]
Isotropy and the profile & \lean{isIsotropic_iff_eq_profile}\\\addlinespace[2pt]
\Cref{lem:profile} & \lean{isProfile_profile}\\\addlinespace[2pt]
$F(A)=\abs A^p\,h(\det A/\abs A^2)$ & \lean{apply_eq_opNorm_rpow_mul},
  \lean{IsProfile.eq_rpow_mul}\\\addlinespace[2pt]
\Cref{lem:supports} & \lean{HasPolyaffineSupportAt.quasiconvexAtOn},
  \lean{polyconvex_of_hasPolyaffineSupportAt}, \lean{Polyconvex.hasPolyaffineSupportAt},
  \lean{Polyconvex.quasiconvex}, \lean{Polyconvex.continuous}, \lean{Polyconvex.rankOneConvex},
  \lean{quasiconvexAtOn_of_le_of_eq}\\\addlinespace[2pt]
\Cref{lem:lift} & \lean{hasPolyaffineSupportAt_of_profile}\\[4pt]
\multicolumn{2}{@{}l}{\emph{\Cref{sec:roc,app:roc}}}\\*[2pt]
\Cref{lem:basic} & \lean{IsProfile.convexOn_slice}, \lean{IsProfile.le_of_sq_sub_sq_eq}
  \\\addlinespace[2pt]
\Cref{def:minorant} & \lean{burkholderProfile}, \lean{burkholderType_eq_profile},
  \lean{burkholder_eq_burkholderType}, \lean{IsProfile.burkholderMinorant_iff_slice}
  \\\addlinespace[2pt]
\Cref{lem:minorant} & \lean{IsProfile.exists_burkholderMinorant_eq}\\\addlinespace[2pt]
\Cref{lem:convex} & \lean{ConvexOn.add_mul_le_of_hasDerivAt_le}\\\addlinespace[2pt]
\Cref{prop:tangent} & \lean{IsProfile.tangent}\\\addlinespace[2pt]
\Cref{cor:BE} & \lean{IsProfile.baker_ericksen}\\\addlinespace[2pt]
\Cref{lem:ends} & \lean{IsProfile.slice_ge_of_one}, \lean{IsProfile.slice_ge_of_neg_one}\\[4pt]
\multicolumn{2}{@{}l}{\emph{\Cref{sec:pge2,app:pge2}}}\\*[2pt]
\Cref{lem:burkholder} & \lean{burkholder_support_le}\\\addlinespace[2pt]
\eqref{eq:caseIII} & \lean{BurkholderMinorant.support}, \lean{BurkholderMinorant.support_two},
  \lean{support_slice_eq}\\\addlinespace[2pt]
\Cref{lem:caseI-algebra} & \lean{caseOne_slope_pos}, \lean{caseOne_const_neg},
  \lean{caseOne_relations}\\\addlinespace[2pt]
\Cref{lem:mu} & \lean{caseOne_slope_bounds}\\\addlinespace[2pt]
\Cref{lem:line} & \lean{IsProfile.support_of_stationary}\\\addlinespace[2pt]
\Cref{lem:psi} & \lean{exists_stationary_ge}\\\addlinespace[2pt]
\Cref{prop:caseI} & \lean{IsProfile.support_of_neg}, \lean{support_of_neg_eq}\\\addlinespace[2pt]
Case II & \lean{IsProfile.support_of_neg_aux'}, \lean{support_of_neg_eq'}\\\addlinespace[2pt]
\Cref{prop:clamp} & \lean{IsProfile.support_multiplier}, \lean{support_multiplier_eq}
  \\\addlinespace[2pt]
\eqref{eq:end-support} & \lean{IsProfile.support_one_zero}, \lean{support_one_zero_eq},
  \lean{IsProfile.support_zero_one}\\[4pt]
\multicolumn{2}{@{}l}{\emph{\Cref{sec:plt2,app:plt2}}}\\*[2pt]
\Cref{prop:low} & \lean{IsProfile.exists_burkholderMinorant_of_lt_two}\\\addlinespace[2pt]
\Cref{lem:pone} & \lean{IsProfile.nonneg_of_one}\\\addlinespace[2pt]
\Cref{lem:sublinear} & \lean{ConvexOn.eq_of_tendsto_div_zero}\\[4pt]
\multicolumn{2}{@{}l}{\emph{\Cref{sec:signed,app:signed}}}\\*[2pt]
The function $h$ & \lean{slice}, \lean{IsProfile.eq_rpow_mul_slice},
  \lean{IsProfile.convexOn_slice_Icc}, \lean{IsProfile.continuousOn_slice}\\\addlinespace[2pt]
\Cref{lem:orient} & \lean{IsProfile.slice_zero_nonneg}, \lean{IsProfile.rpow_mul_slice_neg_one_le},
  \lean{IsProfile.rpow_mul_slice_one_le}, \lean{IsProfile.slice_neg_one_neg_or_slice_one_neg},
  \lean{IsProfile.slope_ge}, \lean{IsProfile.slope_le}, \lean{IsProfile.slope_pos}
  \\\addlinespace[2pt]
\Cref{def:hard} & \lean{IsHardBase}\\\addlinespace[2pt]
\Cref{lem:hard} & \lean{IsProfile.slice_neg_nonneg}, \lean{IsHardBase.base_pos},
  \lean{IsHardBase.slope_pos}, \lean{IsHardBase.e_mul_slice_pos}\\\addlinespace[2pt]
\Cref{lem:easy} & \lean{burkholderTypeCoord_eq_burkholder_add},
  \lean{quasiconvex_burkholderTypeCoord_of_two_lt},
  \lean{quasiconvex_burkholderTypeCoord_of_lt_two}\\\addlinespace[2pt]
\eqref{eq:restricted} & \lean{supportDet}, \lean{supportMod}, \lean{supportConst},
  \lean{restrictedSupport}\\\addlinespace[2pt]
\Cref{lem:rs-identities} & \lean{restrictedSupport_base}, \lean{restrictedSupport_of_add_eq_one},
  \lean{supportDet_sub}, \lean{supportConst_add}, \lean{restrictedSupport_rpow_eq},
  \lean{IsHardBase.supportDet_pos}, \lean{IsHardBase.supportMod_nonneg},
  \lean{IsHardBase.supportConst_neg}\\\addlinespace[2pt]
\Cref{prop:restricted} & \lean{IsHardBase.restrictedSupport_le_of_lt},
  \lean{IsHardBase.restrictedSupport_le_of_nonneg},
  \lean{IsHardBase.restrictedSupport_le_of_add_le_one}, \lean{IsProfile.support_of_lt},
  \lean{IsProfile.support_of_neg_of_nonneg}, \lean{IsHardBase.restrictedSupport_le_of_forall_le}
  \\\addlinespace[2pt]
\Cref{prop:monotone} & \lean{IsHardBase.scaled_monotone}, \lean{IsHardBase.scaled_le},
  \lean{IsHardBase.scaled_eq}\\\addlinespace[2pt]
\Cref{prop:radial} & \lean{HomogeneousIsotropicChallenge.proposition6_8};
  \lean{quasiconvex_radial}, \lean{radial}, \lean{radialProfile}\\\addlinespace[2pt]
\Cref{lem:radial} & \lean{radial_affine}, \lean{radial_hinge}, \lean{radialKernel},
  \lean{radialKernel_one}, \lean{radialKernel_neg_one}, \lean{quasiconvex_radialKernel},
  \lean{abs_radial_sub_le}\\\addlinespace[2pt]
\Cref{lem:hinge} & \lean{ConvexOn.exists_hinge_approx}\\\addlinespace[2pt]
\Cref{cor:burkholder} & \lean{radial_rpow_eq_burkholder},
  \lean{radial_rpow_eq_burkholder_reflection_mul}, \lean{radial_neg_rpow_eq_burkholder},
  \lean{radial_neg_rpow_eq_burkholder_reflection_mul}, \lean{quasiconvex_burkholder_of_one_lt},
  \lean{quasiconvex_burkholder_reflection_mul}, \lean{quasiconvex_burkholder_of_two_le},
  \lean{quasiconvex_burkholder_of_le_two}; the identities \lean{integral_max_rpow},
  \lean{radial_rpow}, \lean{radial_neg_rpow}, \lean{burkholder_eq_burkholderTypeCoord},
  \lean{burkholderTypeCoord_reflection_mul}\\\addlinespace[2pt]
\Cref{def:path} & \lean{IsPath}\\\addlinespace[2pt]
\Cref{prop:path} & \lean{IsHardBase.exists_isPath}\\\addlinespace[2pt]
$M$, $f$, $N$ & \lean{pathSlope}, \lean{pathDet}, \lean{pathEnvelope}\\\addlinespace[2pt]
\Cref{lem:path} & \lean{IsPath.hasDerivAt_pathSlope}, \lean{IsPath.monotoneOn_pathDet},
  \lean{IsPath.pathSlope_add_le}, \lean{pathSlope_deriv_identity},
  \lean{scaled_monotone_identity}\\\addlinespace[2pt]
\Cref{lem:envelope} & \lean{IsPath.convexOn_pathEnvelope}, \lean{IsPath.pathEnvelope_scale},
  \lean{IsPath.pathEnvelope_of_le_one}, \lean{IsPath.scale_mul_pathSlope_sub_integral},
  \lean{scale_rpow_mul_deriv_identity}\\\addlinespace[2pt]
\Cref{prop:minorant} & \lean{IsPath.exists_convex_minorant}\\\addlinespace[2pt]
\Cref{cor:minorant} & \lean{IsHardBase.exists_convex_minorant}\\\addlinespace[2pt]
\Cref{lem:reflect} & \lean{isIsotropic_comp_reflection_mul},
  \lean{IsIsotropic.apply_reflection_mul}, \lean{IsPosHomogeneous.comp_reflection_mul},
  \lean{rankOneConvex_comp_reflection_mul}, \lean{slice_profile_comp_reflection_mul_neg_one},
  \lean{quasiconvex_comp_reflection_mul}\\\addlinespace[2pt]
\Cref{lem:perturb} & \lean{coordPowSum}, \lean{isIsotropic_coordPowSum},
  \lean{isPosHomogeneous_coordPowSum}, \lean{profile_coordPowSum},
  \lean{rankOneConvex_coordPowSum}, \lean{strictConvexOn_add_mul_rpow}\\\addlinespace[2pt]
\Cref{prop:interior} & \lean{quasiconvexAtOn_of_strictConvexOn_slice},
  \lean{quasiconvexAtOn_of_isHardBase}\\\addlinespace[2pt]
\Cref{lem:dense} & \lean{quasiconvex_of_forall_ne_zero}\\\addlinespace[2pt]
\Cref{prop:oriented} & \lean{quasiconvex_of_slice_neg_one_neg},
  \lean{quasiconvex_of_strictConvexOn_slice}\\\addlinespace[2pt]
\Cref{lem:qc-operations} & \lean{quasiconvex_comp_smul}, \lean{quasiconvex_const_add},
  \lean{quasiconvex_finset_sum}, \lean{quasiconvex_add_mul},
  \lean{quasiconvex_norm_add_norm_add_det}\\\addlinespace[2pt]
\Cref{lem:psi-low} & \lean{exists_stationary_ge_of_lt_two}\\\addlinespace[2pt]
\Cref{lem:decreasing} & \lean{add_add_le_of_forall_quadratic_nonpos}\\\addlinespace[2pt]
\Cref{lem:caseII-algebra} & \lean{caseTwo_relations}, \lean{caseTwo_const_neg},
  \lean{caseTwo_slope_bounds}\\\addlinespace[2pt]
\Cref{lem:line-low} & \lean{IsProfile.support_of_stationary_of_lt_two}\\\addlinespace[2pt]
\Cref{lem:legendre} & \lean{StrictConvexOn.exists_legendre}\\\addlinespace[2pt]
\Cref{lem:quadrature} & \lean{intervalIntegral.exists_exp_integral_eq},
  \lean{intervalIntegral.hasDerivAt_exp_integral}\\\addlinespace[2pt]
\Cref{lem:self-improving} & \lean{monotoneOn_of_forall_sub_le_half}\\
\end{longtable}
}


\section*{Acknowledgments and the use of AI}

The results of this paper were found and written with AI systems. ChatGPT Chat~6 Extra High
(OpenAI) and Claude Fable~5.2 (Anthropic) worked on the research that led to the argument for
nonnegative integrands; it is recorded in the first pull request of the
repository~\cite{Companion}. Claude Fable~5.2 wrote that argument as research notes, then the
Lean formalization, in which several steps take shorter routes than the notes, and then the
manuscript from the formalization. The argument for integrands that change sign is due to
ChatGPT Chat~6 Extra High, and is recorded in the same pull request. Claude Fable~5.2
reorganized it so that it needs no theory of ordinary differential equations and no
measure-valued weight, and formalized it; Claude Opus~5.5 (Anthropic) revised the manuscript for
it. Agents running Claude Opus~5.5 wrote the formalization~\cite{IwaniecLean} of~\cite{SharpHI}.
No person has checked the proofs line by line. Lean's kernel checks the formal proofs of the
statements listed in \cref{app:formal-statements}.


\bibliographystyle{amsplain}
\bibliography{refs}

\end{document}
